\documentclass[11pt,a4paper]{article}

% --- Encoding ---
\usepackage[utf8]{inputenc}
\usepackage[T1]{fontenc}
\usepackage[spanish,es-noshorthands]{babel}
\usepackage{lmodern}

% --- Layout ---
\usepackage[margin=2.4cm]{geometry}
\usepackage{parskip}
\setlength{\parindent}{0pt}
\setlength{\parskip}{0.55em plus 0.2em minus 0.1em}

% --- Tables and graphics ---
\usepackage{booktabs}
\usepackage{array}
\usepackage{enumitem}
\usepackage{xcolor}
\usepackage{graphicx}
\usepackage{tikz}
\usetikzlibrary{positioning,shapes.geometric,arrows.meta,fit,backgrounds,calc,decorations.pathreplacing}
\usepackage{caption}
\usepackage{subcaption}
\usepackage{amsmath,amssymb}
\usepackage{longtable}

% --- Soporte para caracteres Unicode ---
\usepackage{newunicodechar}
\newunicodechar{→}{\ensuremath{\rightarrow}}
\newunicodechar{≥}{\ensuremath{\geq}}
\newunicodechar{≤}{\ensuremath{\leq}}
\newunicodechar{×}{\ensuremath{\times}}
\newunicodechar{÷}{\ensuremath{\div}}
\newunicodechar{≈}{\ensuremath{\approx}}
\newunicodechar{⌈}{\ensuremath{\lceil}}
\newunicodechar{⌉}{\ensuremath{\rceil}}
\newunicodechar{✓}{\ensuremath{\checkmark}}
\newunicodechar{≠}{\ensuremath{\neq}}
\newunicodechar{∨}{\ensuremath{\vee}}
\newunicodechar{∧}{\ensuremath{\wedge}}
\newunicodechar{∀}{\ensuremath{\forall}}
\newunicodechar{—}{---}
\newunicodechar{–}{--}
\newunicodechar{←}{\ensuremath{\leftarrow}}
\newunicodechar{↑}{\ensuremath{\uparrow}}
\newunicodechar{↓}{\ensuremath{\downarrow}}
\newunicodechar{∞}{\ensuremath{\infty}}
\newunicodechar{²}{\ensuremath{^{2}}}
\newunicodechar{³}{\ensuremath{^{3}}}
\newunicodechar{⁰}{\ensuremath{^{0}}}
\newunicodechar{¹}{\ensuremath{^{1}}}
\newunicodechar{⁴}{\ensuremath{^{4}}}
\newunicodechar{⁵}{\ensuremath{^{5}}}
\newunicodechar{⁶}{\ensuremath{^{6}}}
\newunicodechar{⁷}{\ensuremath{^{7}}}
\newunicodechar{⁸}{\ensuremath{^{8}}}
\newunicodechar{⁹}{\ensuremath{^{9}}}
\newunicodechar{…}{\ldots}

% --- Code listings ---
\usepackage{listings}
\definecolor{codegray}{rgb}{0.96,0.96,0.96}
\definecolor{codeborder}{rgb}{0.85,0.85,0.85}
\definecolor{codecomment}{rgb}{0.30,0.55,0.30}
\definecolor{codekeyword}{rgb}{0.10,0.30,0.70}
\definecolor{codestring}{rgb}{0.65,0.15,0.10}

\lstdefinestyle{cstyle}{
    language=C,
    backgroundcolor=\color{codegray},
    basicstyle=\ttfamily\footnotesize,
    breaklines=true,
    frame=single,
    rulecolor=\color{codeborder},
    showstringspaces=false,
    keywordstyle=\color{codekeyword}\bfseries,
    commentstyle=\color{codecomment}\itshape,
    stringstyle=\color{codestring},
    aboveskip=1em,
    belowskip=1em,
    columns=fullflexible,
    morekeywords={pte,vaddr,paddr,uint,uint8,uint32,uint64,size_t,address,bool,spinlock,task,trap_frame,context,cpu_state,file,link},
}

% RISC-V assembly style
\lstdefinelanguage{riscv}{
  morekeywords={addi,add,sub,mul,div,la,li,lw,sw,ld,sd,beq,bne,blt,bge,bltu,bgeu,j,jal,jalr,ret,call,mv,mret,sret,ecall,ebreak,wfi,csrr,csrw,csrs,csrc,csrrw,csrrs,csrrc,csrci,csrsi,fence,sfence,amoswap,unimp},
  sensitive=true,
  morecomment=[l]{\#},
  morestring=[b]",
}
\lstdefinestyle{asmstyle}{
    language=riscv,
    backgroundcolor=\color{codegray},
    basicstyle=\ttfamily\footnotesize,
    breaklines=true,
    frame=single,
    rulecolor=\color{codeborder},
    showstringspaces=false,
    keywordstyle=\color{codekeyword}\bfseries,
    commentstyle=\color{codecomment}\itshape,
    stringstyle=\color{codestring},
    aboveskip=1em,
    belowskip=1em,
    columns=fullflexible,
}

\lstdefinestyle{shellstyle}{
    language=bash,
    backgroundcolor=\color{codegray},
    basicstyle=\ttfamily\small,
    breaklines=true,
    frame=single,
    rulecolor=\color{codeborder},
    showstringspaces=false,
    keywordstyle=\color{codekeyword}\bfseries,
    commentstyle=\color{codecomment}\itshape,
    stringstyle=\color{codestring},
    aboveskip=1em,
    belowskip=1em,
    columns=fullflexible,
    upquote=true,
}

\lstdefinestyle{textstyle}{
    backgroundcolor=\color{codegray},
    basicstyle=\ttfamily\footnotesize,
    breaklines=true,
    frame=single,
    rulecolor=\color{codeborder},
    showstringspaces=false,
    aboveskip=1em,
    belowskip=1em,
    columns=fullflexible,
}

\lstset{style=textstyle}

% --- Boxes ---
\usepackage[most]{tcolorbox}
\newtcolorbox{idea}[1][]{
    colback=orange!5!white, colframe=orange!60!black,
    fonttitle=\bfseries, title=Idea intuitiva,
    sharp corners=southwest, rounded corners=northwest, arc=2pt,
    boxrule=0.6pt, left=8pt, right=8pt, top=4pt, bottom=4pt, #1
}
\newtcolorbox{importante}[1][]{
    colback=red!5!white, colframe=red!50!black,
    fonttitle=\bfseries, title=Importante,
    sharp corners=southwest, rounded corners=northwest, arc=2pt,
    boxrule=0.6pt, left=8pt, right=8pt, top=4pt, bottom=4pt, #1
}
\newtcolorbox{respuesta}[1][]{
    colback=green!5!white, colframe=green!40!black,
    fonttitle=\bfseries, title=Respuesta,
    sharp corners=southwest, rounded corners=northwest, arc=2pt,
    boxrule=0.6pt, left=8pt, right=8pt, top=4pt, bottom=4pt, #1
}
\newtcolorbox{walkthrough}[1][]{
    colback=blue!3!white, colframe=blue!50!black,
    fonttitle=\bfseries, title=Recorrido del código,
    sharp corners=southwest, rounded corners=northwest, arc=2pt,
    boxrule=0.6pt, left=8pt, right=8pt, top=4pt, bottom=4pt, #1
}

\usepackage[colorlinks=true,linkcolor=blue!60!black,urlcolor=blue!50!black,citecolor=blue!60!black]{hyperref}

% --- Color palette ---
\definecolor{c1}{HTML}{4C72B0}
\definecolor{c2}{HTML}{DD8452}
\definecolor{c3}{HTML}{55A868}
\definecolor{c4}{HTML}{8172B2}
\definecolor{cred}{HTML}{C44E52}
\definecolor{cfree}{HTML}{8C8C8C}

\title{
    \vspace{-2.5em}
    \textbf{EDOS --- Resumen de estudio} \\[0.3em]
    \large Para el examen práctico y la defensa de implementación \\
    \small Basado en el código entregado en \texttt{08-process/}
}
\author{
    Tomás Rodeghiero \\
    \small Sistemas Operativos --- UNRC
}
\date{}

\begin{document}

\maketitle

\begin{abstract}
\noindent
Este documento es la fuente única para preparar el examen práctico
y la defensa oral de implementación sobre EDOS, el SO educativo
basado en RISC-V utilizado por la cátedra de Sistemas Operativos de
la UNRC. La fuente de verdad es el código entregado en
\texttt{trabajos-entregables/08-process/}; todas las citas a
archivo/línea se refieren a ese árbol.

El documento se organiza en cinco partes: (I) visión global de EDOS,
arquitectura RISC-V y plataforma; (II) recorrido subsistema por
subsistema del kernel con \emph{walkthrough} del código real; (III)
tabla comparativa contra xv6 con la justificación pedagógica de
cada diferencia; (IV) banco de preguntas razonables que un docente
puede hacer en la defensa, con respuesta completa; (V)
\emph{cheat sheet} compacto de registros, formato de PTE, secuencia
de boot, primer \texttt{fork} y comandos útiles de QEMU + GDB.

\begin{importante}
\textbf{Aclaraciones técnicas sobre EDOS.} La consigna mencionaba
``RISC-V 64 (rv64)'' y ``paginación Sv39''. El código entregado
\textbf{no es así}; este documento describe lo que hay en el código:
\begin{itemize}[nosep,leftmargin=*]
    \item EDOS es \textbf{RISC-V 32 bits} (\texttt{-march=rv32ima
        -mabi=ilp32}, \texttt{qemu-system-riscv32}). El ensamblador
        usa \texttt{lw}/\texttt{sw}; los registros guardados ocupan
        4 bytes cada uno.
    \item La paginación es \textbf{Sv32} (dos niveles: PD 10 b +
        PT 10 b + offset 12 b $=$ 32). PTE de 32 bits con PPN de
        22 b y 10 b de flags (\texttt{arch.h:160-213}).
    \item EDOS \textbf{no usa una página \emph{trampoline} separada}
        como xv6. Los handlers \texttt{s\_trap} y \texttt{u\_trap}
        están en el segmento \texttt{.text} del kernel, mapeado en
        cada \texttt{pgtbl} de proceso porque \texttt{exec} copia
        las entradas del \texttt{kernel\_pgtbl} en el inicio del
        nuevo page directory (\texttt{task.c:386}).
    \item EDOS sólo define las flags \textbf{V, R, W, X, U}. No
        expone bits A, D ni G en el PTE (\texttt{arch.h:175-179}).
\end{itemize}
Estas diferencias son centrales para la defensa: si te preguntan
``Sv39 vs Sv32'', tenés que poder explicar por qué EDOS adoptó
Sv32.
\end{importante}
\end{abstract}

\tableofcontents
\newpage

% =====================================================================
\part{Visión global de EDOS}
% =====================================================================

\section{Qué es EDOS y para qué sirve}

EDOS (\emph{Educational Operating System}) es un kernel mínimo
escrito en C y RISC-V assembler, mantenido por la cátedra de
Sistemas Operativos de la UNRC y heredado en espíritu de xv6 (MIT).
Su objetivo no es competir con un SO real: es \textbf{transparentar
los mecanismos del SO}. Cada subsistema (gestor de memoria,
scheduler, trap handler, sistema de archivos, syscalls) está
implementado en la cantidad mínima de líneas que demuestren la idea
con claridad, y vive en archivos pequeños y aislados que el
estudiante puede leer, modificar y depurar con \texttt{gdb} en una
tarde.

\subsection*{Objetivos pedagógicos}

\begin{itemize}[leftmargin=*]
    \item \textbf{Hacer visible} cada paso desde que se enciende la
        CPU hasta que un proceso de usuario corre código.
    \item \textbf{Forzar al alumno a tocar} el código del kernel:
        cada taller agrega un subsistema y rompe algo del anterior
        para que sea evidente qué hay que arreglar.
    \item \textbf{Trabajar con RISC-V}, una arquitectura abierta,
        con manuales accesibles y un conjunto de instrucciones lo
        suficientemente simple para que el ensamblador no oscurezca
        las ideas.
    \item \textbf{Apoyarse en xv6}, que comparte mucha de la filosofía
        pero ya en RISC-V; EDOS se permite simplificar a propósito
        donde xv6 ya es ``demasiado'' para una cursada.
\end{itemize}

\subsection*{Diferencias clave con un SO real}

\begin{itemize}[leftmargin=*,nosep]
    \item No hay un sistema de archivos persistente: \textbf{EFS}
        es \emph{embebido en el kernel} a través de un \texttt{.c}
        que genera \texttt{xxd} (\texttt{efsfiles.c}).
    \item No hay \texttt{fork}/\texttt{exec} desde el espacio de
        usuario: el único \texttt{fork} conceptual es el que el
        kernel hace al crear \texttt{init}.
    \item No hay \texttt{wait}/\texttt{exit} de uso general: lo
        cubre el scheduler matando ZOMBIEs.
    \item No hay manejo elaborado de señales, redes, ni soporte de
        \emph{copy-on-write}.
    \item La planificación es \textbf{round-robin puro}, sin
        prioridades.
\end{itemize}

% ---------------------------------------------------------------------
\section{Arquitectura objetivo: RISC-V 32}

\subsection{Modos de privilegio}

RISC-V define tres modos:

\begin{center}
\small
\begin{tabular}{@{}llp{8.5cm}@{}}
\toprule
Modo & Sigla & Rol \\
\midrule
\textbf{Machine}    & M & Modo más privilegiado. Acceso total al hardware. En EDOS se usa solo durante el \emph{boot} y para reprogramar el timer del \texttt{CLINT}. \\
\textbf{Supervisor} & S & Donde corre el kernel de EDOS después del boot. Maneja paginación (vía \texttt{satp}) y traps. \\
\textbf{User}       & U & Donde corren los procesos. Acceso restringido por la MMU y las flags \texttt{U} de las PTEs. \\
\bottomrule
\end{tabular}
\end{center}

EDOS \emph{delega} casi todas las interrupciones y excepciones de
M-mode hacia S-mode con \texttt{medeleg} y \texttt{mideleg}
(\texttt{arch.s:41-43}). Sólo la interrupción del timer queda en
M-mode (con \texttt{mtvec=m\_trap}), y allí se reprograma el
\texttt{mtimecmp} y se ``rebota'' como interrupción de software en
S-mode levantando \texttt{sip.SSIP}.

\subsection{Registros CSR relevantes para el SO}

\begin{center}
\small
\begin{longtable}{@{}llp{9.5cm}@{}}
\toprule
CSR & Modo & Rol en EDOS \\
\midrule
\texttt{mhartid} & M & ID del hart (CPU lógica). Se lee al inicio y se guarda en \texttt{tp} (\texttt{arch.s:14}). \\
\texttt{mstatus} & M & Estado global. \texttt{MPP}=01 para volver a S-mode con \texttt{mret} (\texttt{arch.s:33-38}). \\
\texttt{mtvec}   & M & Dirección del handler de M-mode (apunta a \texttt{m\_trap}, \texttt{arch.s:58-59}). \\
\texttt{mepc}    & M & PC que va a usarse en el próximo \texttt{mret}. Se setea a \texttt{supervisor} (\texttt{arch.s:69-70}). \\
\texttt{mie}     & M & Habilitación de IRQs de M-mode (\texttt{MTIE=bit 7} para timer). \\
\texttt{mip}     & M & Pendientes de M-mode. \\
\texttt{mtimecmp}& M & Comparador del timer (MMIO en \texttt{CLINT}). Cuando \texttt{mtime} $\geq$ \texttt{mtimecmp} dispara la IRQ. \\
\texttt{medeleg} & M & Delega excepciones a S-mode (EDOS pone todo: \texttt{0xFFFF}). \\
\texttt{mideleg} & M & Delega interrupciones a S-mode (\texttt{0xFFFF}). \\
\texttt{pmpcfg0} & M & Configuración del PMP para que S/U vean toda la RAM (\texttt{arch.s:26-30}). \\
\midrule
\texttt{sstatus} & S & Estado de S-mode. Bit \texttt{SIE} (interrupciones), \texttt{SPP} (modo previo), \texttt{SPIE}. \\
\texttt{stvec}   & S & Dirección del handler de S-mode. EDOS lo cambia entre \texttt{s\_trap} (mientras está en kernel) y \texttt{u\_trap} (antes de volver a user). \\
\texttt{sepc}    & S & PC en el que ocurrió el trap (donde \texttt{sret} retornará). \\
\texttt{scause}  & S & Causa del trap (interrupción vs.\ excepción y código). \\
\texttt{stval}   & S & Dato auxiliar (ej.: dirección que generó un page fault). \\
\texttt{sscratch}& S & ``Bolsillo'' del kernel. En EDOS guarda la base del kernel-stack del proceso, para que \texttt{u\_trap} pueda saltar a él en una instrucción. \\
\texttt{satp}    & S & Page-table address pointer. Bits altos: modo (0 = bare, 1 = Sv32 en rv32). Bits bajos: PPN del directorio. \\
\texttt{sie}     & S & Habilitación de IRQs de S-mode. \\
\texttt{sip}     & S & Pendientes de S-mode. EDOS usa \texttt{SSIP} (software int.) como vehículo para el timer rebotado desde M-mode. \\
\midrule
\multicolumn{3}{@{}l@{}}{\textbf{Registros de propósito general (ABI rv32)}}\\
\midrule
\texttt{ra}     & --- & Return address. \\
\texttt{sp}     & --- & Stack pointer. \\
\texttt{tp}     & --- & ``Thread pointer''. EDOS lo abusa para guardar el \texttt{hartid} (línea \texttt{arch.s:14}: \texttt{csrr tp, mhartid}). \\
\texttt{a0--a7} & --- & Argumentos y resultado de funciones / syscalls. \\
\texttt{s0--s11}& --- & Callee-saved. Son los únicos guardados por \texttt{context\_switch}. \\
\texttt{t0--t6} & --- & Caller-saved. Sólo se salvan en el trap frame. \\
\bottomrule
\end{longtable}
\end{center}

\subsection{Plataforma: QEMU \texttt{virt}}

EDOS corre en la máquina virtual \texttt{virt} de QEMU, lanzada
desde el \texttt{Makefile} con:

\begin{lstlisting}[style=shellstyle]
qemu-system-riscv32 -machine virt -bios none -nographic -smp 1 -kernel kernel
\end{lstlisting}

Mapa de E/S relevante (definido en \texttt{arch.h:18-41}):

\begin{center}
\small
\begin{tabular}{@{}lll@{}}
\toprule
Dispositivo & Base MMIO & Notas \\
\midrule
\texttt{CLINT}   & \texttt{0x02000000} & Timer + IPI por hart. \texttt{mtime} en \texttt{+0xBFF8}, \texttt{mtimecmp[h]} en \texttt{+0x4000+8h}. \\
\texttt{PLIC}    & \texttt{0x0C000000} & Controlador de IRQs externas (UART, virtio). Tamaño 4\,MiB. \\
\texttt{UART}    & \texttt{0x10000000} & NS16550A. Registros HR (rx/tx), LSR, IER (\texttt{console.c:6-12}). \\
\texttt{VIRTIO0} & \texttt{0x10001000} & Disco virtio (reservado; EDOS no lo usa todavía). \\
\midrule
RAM (kernel + apps) & \texttt{0x80000000} & 128 MiB (\texttt{kernel.ld:35}). \\
\bottomrule
\end{tabular}
\end{center}

\textbf{Constantes derivadas:}

\begin{itemize}[leftmargin=*,nosep]
    \item \texttt{T\_INTERVAL = 5\,000\,000}: ticks del \texttt{mtime}
        por cuanto. Aprox.\ medio segundo en QEMU
        (\texttt{arch.h:27}).
    \item \texttt{UART\_IRQ = 10}: número de IRQ del UART en el
        PLIC.
    \item \texttt{NCPU = 4}: cantidad máxima de harts soportada
        (estructuras como \texttt{cpus\_state[]} reservan ese tamaño,
        aunque el \texttt{Makefile} arranca QEMU con \texttt{-smp 1}).
\end{itemize}

% ---------------------------------------------------------------------
\section{Mapa de memoria física}

El linker script \texttt{kernel.ld} fija el layout en RAM:

\begin{center}
\begin{tikzpicture}[
    every node/.style={font=\ttfamily\small},
    sec/.style={draw,minimum width=6.5cm,minimum height=0.75cm,align=center}
]
\node[sec,fill=c1!25] (text)   at (0, 6) {\texttt{.text} (kernel code)};
\node[sec,fill=c1!15,below=0pt of text] (rodata) {\texttt{.rodata}};
\node[sec,fill=c2!25,below=0pt of rodata] (data)  {\texttt{.data}};
\node[sec,fill=c2!15,below=0pt of data]   (bss)   {\texttt{.bss}};
\node[sec,fill=c3!25,below=0pt of bss]    (stk)   {\texttt{\_\_stack0}: 4 stacks de 4\,KiB};
\node[sec,fill=c4!25,below=0pt of stk,minimum height=1.5cm]   (heap) {Pool de páginas libres \\ administrado por \texttt{kalloc}};

\node[anchor=west,xshift=0.5cm] at (text.east)   {\texttt{0x80000000} = \texttt{\_\_kernel\_start}};
\node[anchor=west,xshift=0.5cm] at (rodata.east) {};
\node[anchor=west,xshift=0.5cm] at (data.east)   {\texttt{\_\_text\_end} (alineado a 4\,KiB)};
\node[anchor=west,xshift=0.5cm] at (bss.east)    {\texttt{\_\_bss}, \texttt{\_\_bss\_end}};
\node[anchor=west,xshift=0.5cm] at (stk.east)    {\texttt{\_\_kernel\_end}};
\node[anchor=west,xshift=0.5cm] at (heap.south east) {\texttt{\_\_mem\_end} = base + 128\,MiB};

\end{tikzpicture}
\end{center}

\textbf{Datos importantes:}

\begin{itemize}[leftmargin=*]
    \item Las primeras dos páginas de cada hart son su stack
        inicial. \texttt{arch.s:17-21} asigna a la CPU $h$ el rango
        \texttt{[\_\_stack0 + h$\cdot$4\,KiB, \_\_stack0 + (h+1)$\cdot$4\,KiB)}.
    \item Todo lo que está entre \texttt{\_\_kernel\_end} y
        \texttt{\_\_mem\_end} es el \textbf{pool de páginas} que
        \texttt{init\_kalloc()} encadena en su free-list
        (\texttt{kalloc.c:43-51}).
    \item Las regiones MMIO (\texttt{CLINT}, \texttt{UART},
        \texttt{PLIC}, \texttt{VIRTIO0}) están \emph{fuera} de este
        rango pero se mapean 1:1 en el \texttt{kernel\_pgtbl} para
        que el kernel pueda leerlas/escribirlas con direcciones
        normales (\texttt{arch.c:97-117}).
\end{itemize}

% ---------------------------------------------------------------------
\section{Flujo de arranque end-to-end}

\subsection{Del firmware/OpenSBI al kernel}

QEMU arranca con \texttt{-bios none}, así que \textbf{no hay
OpenSBI} en EDOS: el kernel mismo se carga en \texttt{0x80000000} y
la primera instrucción que se ejecuta es \texttt{boot}
(\texttt{arch.s:11-72}). Es un caso atípico: en hardware real o con
firmware, OpenSBI corre primero en M-mode y delega a S-mode; en
EDOS asumimos esa función.

\subsection{Lo que hace \texttt{boot} antes de cualquier C}

\begin{walkthrough}
\textbf{\texttt{arch.s:11-72}}

\begin{enumerate}[leftmargin=*,nosep,label=\arabic*.]
    \item \texttt{csrr tp, mhartid} — guarda el id del hart en
        \texttt{tp}. Es el único lugar donde se lee \texttt{mhartid}.
    \item \texttt{sp = \_\_stack0 + 4\,KiB$\cdot$(hartid+1)} —
        cada hart toma una página de stack distinta. Sin esto,
        todos escribirían en el mismo stack y el primer
        \texttt{call} corrompería el de los otros.
    \item \texttt{csrw satp, x0} — \textbf{deshabilita la paginación}
        para correr en modo \texttt{bare}. Se va a habilitar después
        en \texttt{enable\_paging} con \texttt{satp = 0x80000000 |
        (kernel\_pgtbl>>12)}.
    \item \texttt{pmpcfg0 = 0x1F, pmpaddr0 = -1} — Physical Memory
        Protection: permite a S y U-mode acceder \emph{a toda} la
        memoria física. Sin esto, las trampas se dispararían apenas
        bajáramos a S-mode.
    \item Cambiar \texttt{mstatus.MPP} de 11 (M) a 01 (S) — el
        \texttt{mret} de más abajo va a volver a S, no a M.
    \item \texttt{medeleg = mideleg = 0xFFFF} — todas las
        excepciones e interrupciones se atienden en S-mode. Es el
        truco que evita escribir handlers de M-mode para cada caso.
    \item \texttt{sstatus.SIE = 1} — habilitar IRQs de S-mode.
    \item \texttt{sie} ← \texttt{SEIE | STIE | SSIE} — habilitar
        las tres clases de IRQs de S-mode (externa, timer, software).
    \item \texttt{call next\_timer\_interrupt} — agendar la primera
        IRQ del timer (programa \texttt{mtimecmp}).
    \item \texttt{mtvec = m\_trap} — handler de M-mode (sólo lo
        usa el timer).
    \item Habilitar \texttt{mstatus.MIE} y \texttt{mie.MTIE}.
    \item \texttt{mepc = supervisor; mret} — salta a la etiqueta
        \texttt{supervisor} \emph{ya en S-mode}.
\end{enumerate}
\end{walkthrough}

\subsection{De \texttt{supervisor:} a \texttt{kernel\_main}}

La etiqueta \texttt{supervisor} (\texttt{arch.s:77-83}) corre en
S-mode con el mismo stack inicial. Lo único que hace es:

\begin{lstlisting}[style=asmstyle]
supervisor:
    la      t1, s_trap
    csrw    stvec, t1     # de aca en adelante, traps S-mode -> s_trap
    call    kernel_main   # nunca retorna; arranca el scheduler
\end{lstlisting}

\subsection{Lo que hace \texttt{kernel\_main} (\texttt{kmain.c:10-32})}

\begin{walkthrough}
\textbf{\texttt{kmain.c:10-32}}

\begin{enumerate}[leftmargin=*,nosep,label=\arabic*.]
    \item Sólo el hart 0:
        \begin{enumerate}[nosep,label=\alph*)]
            \item \texttt{init\_kalloc()} — encadena la free-list
                de páginas físicas (\texttt{kalloc.c:43-51}).
            \item \texttt{init\_vm()} — crea el \texttt{kernel\_pgtbl}
                y mapea kernel + MMIO 1:1 (\texttt{vm.c:6-11},
                \texttt{arch.c:97-117}).
            \item \texttt{create\_process("init")} — el primer
                proceso de usuario nace acá. Lo veremos en detalle.
            \item \texttt{ready = 1; \_\_sync\_synchronize()} —
                publica el dato a los otros harts.
        \end{enumerate}
    \item Todos los harts: \texttt{while(!ready) ;} y luego
        \texttt{enable\_paging()} (activa Sv32 en este hart) y
        \texttt{scheduler()} (no retorna).
\end{enumerate}
\end{walkthrough}

\begin{idea}
La sincronización entre harts es deliberadamente cruda: una variable
\texttt{volatile int ready} y un \texttt{\_\_sync\_synchronize()}.
Es válido porque sólo el hart 0 escribe y todos los demás giran
hasta verla en \texttt{1}. En xv6 hay un mecanismo equivalente
(la espera en \texttt{started}).
\end{idea}

% =====================================================================
\part{Subsistemas del kernel}
% =====================================================================

% ---------------------------------------------------------------------
\section{Bootstrap y SMP (arch.s, arch.c, kmain.c)}

\subsection*{Qué resuelve}

Llevar a la CPU desde el reset (en M-mode, sin paginación, con
\texttt{tp} indefinido) hasta un estado en el cual:

\begin{itemize}[nosep,leftmargin=*]
    \item Estamos en S-mode con interrupciones habilitadas.
    \item Cada hart tiene su propio stack.
    \item La paginación está montada con el \texttt{kernel\_pgtbl}
        (en cada hart).
    \item El \texttt{scheduler} de cada hart elegirá una tarea
        \texttt{RUNNABLE} y le dará CPU.
\end{itemize}

\subsection*{Estructuras y constantes clave}

\begin{itemize}[leftmargin=*,nosep]
    \item \texttt{NCPU = 4} (\texttt{arch.h:13}).
    \item \texttt{volatile int ready} en \texttt{kmain.c:8} —
        bandera de sincronización inter-hart.
\end{itemize}

\subsection*{Funciones principales}

\begin{itemize}[leftmargin=*]
    \item \texttt{boot} en \texttt{arch.s:11-72} — handler de
        firmware. Configura PMP, MPP, deleg, IRQs y salta a
        \texttt{supervisor:}.
    \item \texttt{supervisor:} en \texttt{arch.s:77-83} — instala
        \texttt{s\_trap} y llama \texttt{kernel\_main}.
    \item \texttt{kernel\_main(void)} en \texttt{kmain.c:10-32} —
        inicialización (sólo hart 0) y entrega al \texttt{scheduler}.
    \item \texttt{init\_kalloc()}, \texttt{init\_vm()},
        \texttt{create\_process("init")}, \texttt{enable\_paging()},
        \texttt{scheduler()} (cada uno en su archivo).
\end{itemize}

\subsection*{Invariantes}

\begin{itemize}[nosep,leftmargin=*]
    \item Después de \texttt{boot}: cada hart está en S-mode, con
        su propio stack y con \texttt{stvec=s\_trap}.
    \item Sólo el hart 0 ejecuta \texttt{init\_kalloc/init\_vm}.
        Los demás esperan.
    \item Cuando un hart pasa el \texttt{while(!ready)}, el
        \texttt{kernel\_pgtbl} ya existe y es seguro encender
        \texttt{satp}.
\end{itemize}

\subsection*{Errores frecuentes al modificarlo}

\begin{itemize}[nosep,leftmargin=*]
    \item Si \texttt{ready} se asigna sin \texttt{\_\_sync\_synchronize()},
        las escrituras al \texttt{kernel\_pgtbl} pueden no ser
        visibles a otros harts cuando salen del busy-wait.
    \item Si todos los harts llaman a \texttt{init\_vm}, se reescribe
        múltiples veces el page table raíz y se generan dobles
        liberaciones de páginas.
    \item Si se omite el ajuste de \texttt{sp} por hartid, todos
        comparten stack.
\end{itemize}

% ---------------------------------------------------------------------
\section{Consola y \texttt{kprintf} (console.c, klib.c)}

\subsection*{Qué resuelve}

Permitir al kernel imprimir mensajes y leer caracteres por la UART,
con un \texttt{printf} reducido (sólo \%s, \%d, \%x, \%\%) protegido
por un spinlock para no entrelazar la salida en SMP.

\subsection*{Funciones clave}

\begin{itemize}[leftmargin=*]
    \item \texttt{console\_putc(char ch)} en \texttt{console.c:24-30}
        — espera con polling a que \texttt{UART\_LSR.TX\_IDLE} esté
        seteado y escribe el carácter en \texttt{UART\_HR}.
    \item \texttt{console\_puts(const char *s)} en
        \texttt{console.c:33-39} — toma \texttt{console\_lock},
        recorre la cadena llamando a \texttt{console\_putc}, libera.
    \item \texttt{console\_read\_char()} en \texttt{console.c:44-52}
        — chequea \texttt{LSR.RX\_READY}; si hay byte, lo lee. Si
        no, devuelve 0. (\emph{No bloquea}; es un polling oportunista.)
    \item \texttt{printf(fmt, \ldots)} en \texttt{klib.c:52-112} —
        adquiere el lock una sola vez para toda la línea, formatea
        y llama \texttt{putchar} (alias de \texttt{console\_putc}).
\end{itemize}

\subsection*{Walkthrough: cómo se imprime un \%x}

\begin{walkthrough}
\textbf{\texttt{klib.c:94-99}}

\begin{lstlisting}[style=cstyle]
case 'x': {
    int value = va_arg(vargs, int);
    for (int i = 7; i >= 0; i--) {
        int nibble = (value >> (i * 4)) & 0xf;
        putchar("0123456789abcdef"[nibble]);
    }
}
\end{lstlisting}

Es la implementación más simple posible: imprime los 8 nibbles del
entero en orden, sin omitir ceros a la izquierda. Es por eso que
todos los logs de EDOS tienen el formato \texttt{0x80004000}.
\end{walkthrough}

\subsection*{Polling vs.\ interrupciones}

EDOS usa \textbf{polling} para el TX (esperamos en un \texttt{while}
a que el bit \texttt{TX\_IDLE} esté en 1). Para el RX, hay un
\emph{hook} de IRQ del UART (\texttt{UART\_IRQ=10} en el PLIC) pero
en la versión 08-process no se activa el path: \texttt{console.c}
no llama a \texttt{console\_interrupt\_handler}, y
\texttt{init\_external\_interrupts/init\_external\_irqs\_in\_cpu}
están comentados en \texttt{kmain.c:17,29}.

\begin{idea}
\textbf{¿Por qué polling en TX?} Porque el UART de QEMU vacía
inmediatamente: la espera nunca dura más de un par de ciclos. En
hardware real (a 115\,200 bps) esto sería inaceptable; ahí se
usaría una cola circular + IRQ TX.
\end{idea}

\subsection*{Por qué el lock es \emph{caller-side}}

\texttt{console\_putc} \textbf{no} toma el lock; sólo
\texttt{console\_puts} y \texttt{printf} lo hacen. Esto permite que
el \texttt{printf} adquiera el lock una vez por línea y no caer en
un patrón ``char a char con un lock por carácter'' que haría que
dos \texttt{printf} concurrentes se intercalaran a nivel de byte.

% ---------------------------------------------------------------------
\section{Spinlocks (spinlock.c)}

\subsection*{Qué resuelve}

Exclusión mutua \textbf{en multiprocesador} sobre estructuras del
kernel (tabla de procesos, free-list de páginas, consola, ticks).
Un spinlock funciona porque las CPUs comparten memoria coherente y
porque RISC-V provee instrucciones atómicas (extensión A).

\subsection*{Tipo y API}

\begin{lstlisting}[style=cstyle]
typedef unsigned int spinlock;     // 0 = libre, 1 = tomado

void acquire(spinlock *lk);
void release(spinlock *lk);
\end{lstlisting}

El tipo es un simple entero. Hay dos primitivas auxiliares no
exportadas:

\begin{itemize}[leftmargin=*,nosep]
    \item \texttt{push\_irq\_off(void)} en \texttt{spinlock.c:7-17}
        — deshabilita interrupciones y mantiene un contador de
        anidamiento (\texttt{cpus\_state[cpu].noff}) más el estado
        previo de \texttt{sstatus.SIE}.
    \item \texttt{pop\_irq\_off(void)} en \texttt{spinlock.c:19-25}
        — decrementa \texttt{noff} y, sólo si llega a 0, restaura
        las IRQs si estaban habilitadas antes del primer
        \texttt{push}.
\end{itemize}

\subsection*{Walkthrough: \texttt{acquire}}

\begin{walkthrough}
\textbf{\texttt{spinlock.c:27-44}}

\begin{lstlisting}[style=cstyle]
void acquire(spinlock *lk) {
    push_irq_off();                              // (1)
    while (__sync_lock_test_and_set(lk, 1) != 0) // (2)
        ;
    __sync_synchronize();                        // (3)
}
\end{lstlisting}

\textbf{(1)} Antes de cualquier intento, \emph{deshabilitamos
interrupciones en esta CPU}. ¿Por qué? Porque si un trap
(particularmente el timer) llegara mientras tenemos el lock, podría
dispararse \texttt{yield()} y dejaríamos a otro proceso girando
para siempre sobre nuestro lock.

\textbf{(2)} \texttt{\_\_sync\_lock\_test\_and\_set} compila a
\texttt{amoswap.w.aq a5, a5, (s1)} en RISC-V: una operación
atómica que escribe 1 en \texttt{*lk} y devuelve el valor previo.
Mientras devuelva 1, alguien más lo tenía y giramos.

\textbf{(3)} La barrera de memoria garantiza que las lecturas y
escrituras dentro de la sección crítica no se reordenen
\emph{antes} del lock.
\end{walkthrough}

\subsection*{Walkthrough: \texttt{release}}

\begin{lstlisting}[style=cstyle]
void release(spinlock *lk) {
    __sync_synchronize();        // (1) las escrituras de la SC ya son visibles
    __sync_lock_release(lk);     // (2) amoswap.w zero, zero, (s1)
    pop_irq_off();               // (3) reactivar IRQs si corresponde
}
\end{lstlisting}

\begin{itemize}[nosep,leftmargin=*]
    \item (1) Barrera \emph{antes} de soltar el lock.
    \item (2) Liberación atómica (escribe 0).
    \item (3) Cierra el \texttt{push\_irq\_off} pareado.
\end{itemize}

\subsection*{Invariantes y errores típicos}

\begin{itemize}[leftmargin=*]
    \item \textbf{Anidamiento}: \texttt{acquire/release} pueden
        anidarse arbitrariamente; las IRQs solo se restablecen
        cuando \texttt{noff} vuelve a 0.
    \item \textbf{Releer sin tener}: si liberás un lock que no
        tomaste, \texttt{noff} se descuenta de más y eventualmente
        las IRQs se rehabilitan en una sección crítica equivocada.
    \item \textbf{Olvidarse del \texttt{push\_irq\_off}} (versión
        sin esta protección) lleva a deadlocks clásicos cuando un
        trap se dispara durante la sección crítica y termina
        intentando tomar el mismo lock.
    \item Comparado con un mutex de userland: aquí \textbf{no se
        duerme}. La CPU gira. Por eso sólo conviene para regiones
        muy cortas. Para esperas largas, EDOS usa
        \texttt{suspend/wakeup} (sleep-locks de hecho construidos
        sobre spinlocks).
\end{itemize}

% ---------------------------------------------------------------------
\section{Traps e interrupciones (trap.s, trap.c)}

\subsection*{Qué resuelve}

Recibir cualquier evento que cambia el flujo de ejecución
(interrupciones del timer, IRQs de dispositivos, syscalls,
excepciones como page fault) y dispatcharlo al manejador
correspondiente, dejando la CPU en un estado consistente para
retornar.

\subsection*{Estructura: \texttt{struct trap\_frame}}

Definida en \texttt{arch.h:100-132}, contiene 31 enteros de 32 b
(124 B en total):

\begin{lstlisting}[style=cstyle]
struct trap_frame {
    size_t ra, gp, t0..t6, a0..a7, s0..s11, sp, pc;
} __attribute__((packed));
\end{lstlisting}

\textbf{Detalles importantes:}

\begin{itemize}[leftmargin=*,nosep]
    \item No incluye \texttt{x0} (siempre cero) ni \texttt{tp}
        (lo usamos como hartid permanente).
    \item El campo \texttt{pc} es donde se escribe \texttt{sepc}
        al entrar y desde donde se vuelca antes de \texttt{sret}.
    \item El \texttt{sp} guardado es el del modo \emph{previo}
        (user stack si veníamos de U-mode).
    \item Los punteros a frame de las syscalls usan helpers como
        \texttt{syscall\_number(tf) = tf->a7}.
\end{itemize}

\subsection*{Los tres handlers de bajo nivel}

\begin{enumerate}[leftmargin=*]
    \item \textbf{\texttt{s\_trap}} (\texttt{trap.s:10-82}): traps
        en S-mode. Reserva 29$\times$4 B en el stack actual y
        guarda \texttt{ra, gp, t*, a*, s*} ahí. Llama a
        \texttt{kernel\_trap}. Al volver, restaura todo y hace
        \texttt{sret}.
    \item \textbf{\texttt{u\_trap}} (\texttt{trap.s:89-140}): traps
        en U-mode. Es más complejo porque hay que cambiar de stack:
        \begin{itemize}[nosep]
            \item Intercambia \texttt{sp} y \texttt{sscratch}
                (\texttt{csrrw sp, sscratch, sp}). Tras eso,
                \texttt{sp} apunta al kernel-stack del proceso y
                \texttt{sscratch} al user-stack.
            \item Reserva 31$\times$4 B (un \texttt{trap\_frame}
                completo) y guarda los registros, incluido el
                \texttt{sp} de usuario (leído del \texttt{sscratch})
                y \texttt{sepc} en el campo \texttt{pc} del frame.
            \item Llama a \texttt{user\_trap}.
        \end{itemize}
    \item \textbf{\texttt{m\_trap}} (\texttt{trap.s:200-250}): trap
        de M-mode, sólo para el timer. Reprograma
        \texttt{mtimecmp}, escribe \texttt{sip = 2} (setea
        \texttt{SSIP}, que el S-mode verá como interrupción de
        software) y hace \texttt{mret}.
\end{enumerate}

\subsection*{Diagrama: dónde está cada handler}

\begin{center}
\begin{tikzpicture}[
    every node/.style={font=\ttfamily\small,align=center},
    box/.style={draw,rounded corners,minimum width=3.5cm,minimum height=0.8cm},
    user/.style={box,fill=c1!25},
    sup/.style={box,fill=c3!25},
    mach/.style={box,fill=c2!25},
    arr/.style={->,thick}
]
\node[user] (umode) at (0,4) {Modo Usuario};
\node[sup]  (smode) at (5,4) {Modo Supervisor};
\node[mach] (mmode) at (10,4) {Modo Machine};

\node[box,fill=c1!15] (utrap) at (0,2) {\texttt{u\_trap}\\ \scriptsize en trap.s};
\node[box,fill=c3!15] (strap) at (5,2) {\texttt{s\_trap}\\ \scriptsize en trap.s};
\node[box,fill=c2!15] (mtrap) at (10,2) {\texttt{m\_trap}\\ \scriptsize en trap.s};

\node[box,fill=c1!15] (uth)   at (0,0)  {\texttt{user\_trap}\\ \scriptsize en trap.c};
\node[box,fill=c3!15] (kth)   at (5,0)  {\texttt{kernel\_trap}\\ \scriptsize en trap.c};
\node[box,fill=c2!15] (rtt)   at (10,0) {timer + sip.SSIP\\ \scriptsize delega a S-mode};

\draw[arr] (umode) -- node[right,font=\scriptsize]{trap} (utrap);
\draw[arr] (smode) -- node[right,font=\scriptsize]{trap} (strap);
\draw[arr] (mmode) -- node[right,font=\scriptsize]{timer} (mtrap);
\draw[arr] (utrap) -- (uth);
\draw[arr] (strap) -- (kth);
\draw[arr] (mtrap) -- (rtt);
\end{tikzpicture}
\end{center}

\subsection*{Walkthrough: el timer en U-mode}

\begin{walkthrough}
\textbf{Camino completo del timer cuando un proceso está corriendo}

\begin{enumerate}[leftmargin=*]
    \item El \texttt{mtime} alcanza \texttt{mtimecmp}. El hardware
        levanta una IRQ de M-mode.
    \item Al estar deshabilitada la M-mode preemption con \texttt{wfi},
        la CPU salta a \texttt{mtvec} = \texttt{m\_trap}.
    \item \texttt{m\_trap} reprograma \texttt{mtimecmp} (siguiente
        cuanto) y levanta \texttt{sip.SSIP} (software int. en
        S-mode). \texttt{mret} vuelve al código U-mode interrumpido.
    \item Inmediatamente \emph{la propia} CPU ve la IRQ de software
        en S-mode y atrapa a \texttt{stvec = u\_trap} (porque
        estamos en U).
    \item \texttt{u\_trap} swappea \texttt{sp/sscratch}, llena el
        \texttt{trap\_frame} del proceso, llama \texttt{user\_trap}.
    \item \texttt{user\_trap} (\texttt{trap.c:107-181}) reconoce
        \texttt{cause = TIMER\_INTERRUPT} (\texttt{0x80000001},
        bit de signo prendido $\Rightarrow$ interrupción), llama a
        \texttt{ack\_timer\_interrupt()} (limpia \texttt{sip.SSIP})
        y decrementa \texttt{task->ticks}.
    \item Si llega a 0, llama a \texttt{yield()} (queda
        \texttt{RUNNABLE} y vuelve al scheduler).
    \item Cuando vuelve la CPU al proceso, \texttt{return\_to\_user\_mode()}
        y \texttt{u\_trap\_ret} restauran el frame.
\end{enumerate}
\end{walkthrough}

\subsection*{La función pivote: \texttt{return\_to\_user\_mode}}

\texttt{trap.c:72-101} es la única vía para volver a U-mode. Hace
cinco cosas, todas con interrupciones deshabilitadas:

\begin{enumerate}[leftmargin=*,nosep]
    \item Instalar \texttt{stvec = u\_trap}.
    \item Guardar \texttt{sscratch} = base + 4\,KiB del kernel-stack
        del proceso. \emph{Esa} dirección la usará \texttt{u\_trap}
        para conmutar de pila la próxima vez.
    \item Setear \texttt{sstatus.SPP=0} y \texttt{SPIE=1}, que el
        \texttt{sret} interpretará como ``vamos a U-mode con IRQs
        habilitadas''.
    \item Cargar la page-table del proceso con \texttt{set\_page\_table}
        (que también hace los \texttt{sfence.vma}).
    \item Saltar a \texttt{u\_trap\_ret(task->trap\_frame)} en
        ensamblador, que mueve \texttt{sp = a0}, restaura
        \texttt{sepc} desde el frame, restaura registros y hace
        \texttt{sret}.
\end{enumerate}

% ---------------------------------------------------------------------
\section{Memoria física: \texttt{kalloc} (kalloc.c)}

\subsection*{Qué resuelve}

Asignar y liberar \emph{páginas físicas} de 4\,KiB. Todas las
estructuras del kernel que crecen dinámicamente (page tables, kernel
stacks, code/data/stack de procesos) las piden a \texttt{alloc\_page}.

\subsection*{Estructura}

Una \textbf{free-list} encadenada. Cada página libre lleva en sus
primeros bytes un puntero al ``próximo'' (\texttt{struct link}):

\begin{lstlisting}[style=cstyle]
struct link { struct link *next; };
static struct link *free_list = NULL;
static unsigned int free_pages = 0;
\end{lstlisting}

\begin{idea}
\textbf{Por qué los punteros viven dentro de las páginas libres.}
Como cada página libre tiene 4\,KiB \emph{disponibles}, escribir un
puntero en sus primeros 4 bytes no cuesta memoria adicional. El
truco se usa en xv6, Linux y en cualquier SO con \emph{slab/freelist}
de tamaño fijo.
\end{idea}

\subsection*{API}

\begin{itemize}[leftmargin=*,nosep]
    \item \texttt{void* alloc\_page(void)} en \texttt{kalloc.c:16-24}
        — desencola la cabeza de \texttt{free\_list}, hace
        \texttt{memset(\dots, 0, PAGE\_SIZE)} y la devuelve. Si la
        lista está vacía, devuelve 0.
    \item \texttt{void free\_page(void* pa)} en \texttt{kalloc.c:27-40}
        — verifica alineamiento y rango, encola al inicio de
        \texttt{free\_list}.
    \item \texttt{void init\_kalloc(void)} en \texttt{kalloc.c:43-51}
        — desde \texttt{\_\_mem\_end - PAGE\_SIZE} bajando hasta
        \texttt{\_\_kernel\_end}, llama \texttt{free\_page} en
        cada página. Resultado: la free-list queda con la página
        de menor dirección al frente (LIFO inverso).
\end{itemize}

\subsection*{Invariantes y errores típicos}

\begin{itemize}[leftmargin=*,nosep]
    \item Toda página devuelta por \texttt{alloc\_page} está
        \textbf{en cero}.
    \item \texttt{free\_page} \textbf{aborta con panic} si la
        dirección no está alineada a 4\,KiB o cae fuera del rango
        \texttt{[\_\_kernel\_end, \_\_mem\_end]}.
    \item \textbf{No hay protección de doble \texttt{free}}: si
        liberás dos veces la misma página, queda dos veces en la
        free-list y la próxima asignación devuelve un puntero
        compartido. Ese error no se detecta hasta que se manifiesta
        en un crash de otra estructura.
    \item \texttt{alloc\_page} no toma lock; en SMP debería estar
        protegido. En el código entregado esta carrera está latente,
        pero como \texttt{Makefile:31} fuerza \texttt{-smp 1}, el
        problema no aparece.
\end{itemize}

% ---------------------------------------------------------------------
\section{Memoria virtual: paginación Sv32 (vm.c, vm.h, arch.h)}

\subsection*{Qué resuelve}

Construir y mantener las tablas de páginas que la MMU usa para
traducir direcciones virtuales a físicas. Sin esto, el kernel
funcionaría en modo \emph{bare} (acceso directo a RAM física), no
podríamos darle un espacio de direcciones independiente a cada
proceso ni proteger zonas del kernel del acceso de usuario.

\subsection*{Formato de PTE en Sv32}

EDOS usa el régimen \texttt{Sv32} de RISC-V: tablas de \textbf{dos
niveles}, página de \textbf{4\,KiB}. Una dirección virtual de 32
bits se descompone:

\begin{center}
\begin{tikzpicture}[
    every node/.style={font=\ttfamily\small},
    cell/.style={draw,minimum height=0.7cm,align=center}
]
\node[cell,minimum width=2.8cm,fill=c1!20] (vpn1) at (0,0) {VPN[1]};
\node[cell,minimum width=2.8cm,fill=c2!20,right=0pt of vpn1] (vpn0) {VPN[0]};
\node[cell,minimum width=3.4cm,fill=c3!20,right=0pt of vpn0] (off)  {offset};
\node[font=\scriptsize] at ($(vpn1.north)+(0,0.25)$) {10 bits (idx1)};
\node[font=\scriptsize] at ($(vpn0.north)+(0,0.25)$) {10 bits (idx0)};
\node[font=\scriptsize] at ($(off.north)+(0,0.25)$)  {12 bits};
\node[font=\scriptsize] at ($(vpn1.north west)+(0.1,-0.05)$) {31};
\node[font=\scriptsize] at ($(vpn1.north east)+(-0.15,-0.05)$) {22};
\node[font=\scriptsize] at ($(vpn0.north west)+(0.1,-0.05)$) {21};
\node[font=\scriptsize] at ($(vpn0.north east)+(-0.15,-0.05)$) {12};
\node[font=\scriptsize] at ($(off.north west)+(0.1,-0.05)$)  {11};
\node[font=\scriptsize] at ($(off.north east)+(-0.05,-0.05)$) {0};
\end{tikzpicture}
\end{center}

Una entrada de tabla (PTE) es de \textbf{32 bits} con 22 b de PPN y
10 b de flags:

\begin{center}
\begin{tikzpicture}[
    every node/.style={font=\ttfamily\small},
    cell/.style={draw,minimum height=0.7cm,align=center}
]
\node[cell,minimum width=5.6cm,fill=c1!20] (ppn) at (0,0) {PPN (Physical Page Number)};
\node[cell,minimum width=0.6cm,fill=c3!25,right=0pt of ppn] (u)  {U};
\node[cell,minimum width=0.6cm,fill=c3!25,right=0pt of u]   (x)  {X};
\node[cell,minimum width=0.6cm,fill=c3!25,right=0pt of x]   (w)  {W};
\node[cell,minimum width=0.6cm,fill=c3!25,right=0pt of w]   (r)  {R};
\node[cell,minimum width=0.6cm,fill=c3!25,right=0pt of r]   (v)  {V};
\node[font=\scriptsize] at ($(ppn.north)+(0,0.25)$) {22 bits};
\node[font=\scriptsize] at ($(v.north)+(0,0.25)$)   {1};
\end{tikzpicture}
\end{center}

(EDOS no usa los bits A, D ni G del Sv32 estándar; el RISC-V los
ignora si su hardware no los actualiza, lo cual es lo que pasa en
QEMU.)

\textbf{Flags definidas en \texttt{arch.h:175-179}:}

\begin{itemize}[nosep,leftmargin=*]
    \item \texttt{PAGE\_V (0x1)}: válido. Si está en 0, cualquier
        acceso genera \emph{page fault}.
    \item \texttt{PAGE\_R (0x2)}: lectura permitida.
    \item \texttt{PAGE\_W (0x4)}: escritura permitida.
    \item \texttt{PAGE\_X (0x8)}: ejecución permitida.
    \item \texttt{PAGE\_U (0x10)}: accesible desde U-mode (sin esto,
        el proceso de usuario no puede tocar la página aunque la
        tenga mapeada).
\end{itemize}

\subsection*{Macros de traducción (\texttt{arch.h:188-212})}

\begin{lstlisting}[style=cstyle]
#define page_base_addr(addr) (addr & 0xfffff000)    // round-down a frontera
#define va_offset(va)        (va & 0xfff)           // bits 11..0
#define va_index1(va)        ((va >> 22) & 0x3ff)   // bits 31..22
#define va_index0(va)        ((va >> 12) & 0x3ff)   // bits 21..12

#define pte_ppn(pte)         (pte >> 10)            // PPN en bits 31..10
#define pa2ppn(pa)           ((pa / PAGE_SIZE) << 10)
#define pte_pa(pte)          ((pte >> 10) * PAGE_SIZE)
#define pa2satp(pa)          (0x80000000 | (pa >> 12))
\end{lstlisting}

\textbf{Atención}: el bit alto de \texttt{satp} en Sv32 es el bit
de modo. EDOS hace \texttt{0x80000000 | (pa >> 12)} que en rv32
encodifica ``MODE=1, ASID=0, PPN=pa/4K'' (Sv32 enabled).

\subsection*{Walkthrough: \texttt{get\_pte} y \texttt{map\_page}}

\begin{walkthrough}
\textbf{\texttt{arch.c:48-66} (\texttt{get\_pte}, dos niveles)}

\begin{lstlisting}[style=cstyle]
static pte* get_pte(pte* pgtbl, vaddr va, bool create_leaf) {
    uint i = va_index1(va);
    if ((pgtbl[i] & PAGE_V) == 0) {     // nivel 1 ausente
        if (create_leaf) {              //   y queremos crearlo
            paddr pa = (paddr) alloc_page();
            if (!pa) panic("alloc_page() failed!");
            pgtbl[i] = pa2ppn(pa) | PAGE_V;  // anclamos hoja
        } else return 0;
    }
    pte* pg0 = (pte*) pte_pa(pgtbl[i]); // base de la hoja
    return pg0 ? &pg0[va_index0(va)] : 0;
}
\end{lstlisting}

Dos detalles claves:

\begin{itemize}[nosep,leftmargin=*]
    \item En el nivel 1, el PTE \textbf{sólo} tiene la flag
        \texttt{V} prendida (sin R/W/X). Es la convención del Sv32
        para indicar ``no es una hoja, sigan al próximo nivel''.
    \item La hoja (\texttt{pg0}) es accedida por dirección física,
        porque dentro del kernel todo está mapeado 1:1.
\end{itemize}

\textbf{\texttt{arch.c:69-73} (\texttt{map\_page})}

\begin{lstlisting}[style=cstyle]
void map_page(pte* pgtbl, vaddr va, paddr pa, uint flags) {
    pte* entry = get_pte(pgtbl, va, true);
    *entry = pa2ppn(pa) | flags | PAGE_V;
}
\end{lstlisting}

Una línea para crear la PTE final. Las flags incluyen los
permisos R/W/X/U que pasa el caller.
\end{walkthrough}

\subsection*{Construcción del \texttt{kernel\_pgtbl}}

\begin{walkthrough}
\textbf{\texttt{vm.c:6-11} (\texttt{init\_vm}) y \texttt{arch.c:97-117} (\texttt{map\_kernel\_memory})}

\begin{lstlisting}[style=cstyle]
void init_vm(void) {
    extern pte* kernel_pgtbl;
    kernel_pgtbl = (pte*) alloc_page(); // 1 página = directorio raíz
    map_kernel_memory(kernel_pgtbl);
}

void map_kernel_memory(pte* pgtbl) {
    address ktext_start = (address) __kernel_start;
    uint    ktext_size  = (uint)(__text_end - __kernel_start);
    address kdata_start = (address) (__text_end);
    uint    kdata_size  = (uint) (__mem_end - __text_end);

    map_region(pgtbl, UART,        UART,        PAGE_SIZE, PAGE_R | PAGE_W);
    map_region(pgtbl, VIRTIO0,     VIRTIO0,     PAGE_SIZE, PAGE_R | PAGE_W);
    map_region(pgtbl, PLIC,        PLIC,        PLIC_SIZE, PAGE_R | PAGE_W);
    map_region(pgtbl, ktext_start, ktext_start, ktext_size, PAGE_R | PAGE_X);
    map_region(pgtbl, kdata_start, kdata_start, kdata_size, PAGE_R | PAGE_W);
}
\end{lstlisting}

Cinco regiones \textbf{idénticas en virtual y físico} (1:1):

\begin{enumerate}[nosep,leftmargin=*]
    \item UART (\texttt{0x10000000}): R/W.
    \item VIRTIO0 (\texttt{0x10001000}): R/W.
    \item PLIC (\texttt{0x0c000000}, 4\,MiB): R/W.
    \item Kernel \texttt{.text}: R+X (no W para que un bug del
        kernel no pueda sobrescribir su propio código).
    \item Kernel \texttt{.data} + \texttt{.bss} + stacks + pool de
        páginas: R+W (no X).
\end{enumerate}

Ninguna de estas tiene el bit \texttt{U}, así que los procesos no
pueden tocarlas.
\end{walkthrough}

\subsection*{Cambio de espacio de direcciones (\texttt{arch.h:420-425})}

\begin{lstlisting}[style=cstyle]
static inline void set_page_table(pte* pgtbl) {
    __asm__ __volatile__("sfence.vma");           // invalida TLB anterior
    w_satp(pa2satp((size_t) pgtbl));              // cambia raíz
    __asm__ __volatile__("sfence.vma");           // invalida TLB nueva
}
\end{lstlisting}

Las dos \texttt{sfence.vma} son críticas. Sin ellas, la MMU podría
seguir usando entradas cacheadas que apuntan a la tabla anterior y
toda la traducción quedaría mal.

\subsection*{Resumen de funciones de \texttt{vm.h}}

\begin{itemize}[leftmargin=*,nosep]
    \item \texttt{init\_vm()} — crea \texttt{kernel\_pgtbl}.
    \item \texttt{map\_region(pgtbl, va, pa, size, flags)} — mapea
        un rango (\texttt{vm.c:14-25}).
    \item \texttt{unmap\_region(pgtbl, va, size, free)} — desmapea
        rango (\texttt{vm.c:28-36}).
    \item \texttt{copy\_from\_user(pgtbl, dst, src\_va, n)} y
        \texttt{copy\_to\_user} — traducen la VA de usuario a su PA
        kernel-visible vía \texttt{va2kernel\_address} y hacen
        \texttt{memcpy}.
\end{itemize}

\subsection*{Errores típicos en VM}

\begin{itemize}[leftmargin=*,nosep]
    \item \textbf{Olvidar \texttt{sfence.vma}}: traducción
        impredecible tras cambiar \texttt{satp}.
    \item \textbf{No prender \texttt{PAGE\_U}} en páginas del
        proceso: el kernel las ve OK, pero el proceso obtiene page
        fault al primer acceso.
    \item \textbf{No copiar \texttt{kernel\_pgtbl} en
        \texttt{exec}}: el proceso pierde los mapeos del kernel y,
        cuando entre a un trap, \texttt{s\_trap}/\texttt{u\_trap}
        no estarán en su espacio $\Rightarrow$ panic.
\end{itemize}

% ---------------------------------------------------------------------
\section{Tareas y scheduler (task.c, task.h)}

\subsection*{Qué resuelve}

Multiplexar la CPU entre varias \emph{tareas}. Una \emph{tarea} en
EDOS es la abstracción común que cubre tanto threads de kernel como
procesos de usuario; la diferencia es que un proceso tiene
\texttt{pid > 0} y un \texttt{pgtbl} distinto del \texttt{kernel\_pgtbl}
(\texttt{task.h:67-68}: \texttt{is\_process}).

\subsection*{Estructura \texttt{struct task} (task.h:23-39)}

\begin{lstlisting}[style=cstyle]
struct task {
    uint            tid;             // task id (siempre)
    uint            pid;             // process id (0 si kernel-thread)
    char            name[81];        // nombre amistoso
    int             state;           // UNUSED, CREATED, RUNNABLE, ...
    struct task    *parent;          // padre (NULL para init)
    bool            killed;          // marca de muerte
    int             exit_code;
    struct context  ctx;             // ra, sp, s0..s11 (callee-saved)
    int             ticks;           // cuanto restante (QUANTUM=2 inicial)
    uint8          *kstack;          // base del kernel stack (4 KiB)
    uint64          sleep_ticks;     // contador para sleep(n)
    int             cpu_id;          // CPU donde corre, -1 si no corre
    void           *wait_condition;  // canal de espera (suspend)
    pte            *pgtbl;           // page table del proceso
    spinlock        lock;
};
\end{lstlisting}

\textbf{Tabla global}: \texttt{static struct task tasks[TASK\_MAX]}
con \texttt{TASK\_MAX = 16} (\texttt{task.h:7, task.c:17}).

\subsection*{Máquina de estados}

\begin{center}
\begin{tikzpicture}[
    every node/.style={font=\ttfamily\small},
    state/.style={draw,rounded corners,minimum width=2.6cm,minimum height=0.7cm,align=center},
    arr/.style={->,thick}
]
\node[state,fill=cfree!25] (un) at (0,4)       {UNUSED};
\node[state,fill=c1!25]    (cr) at (4,4)       {CREATED};
\node[state,fill=c2!25]    (rn) at (8,4)       {RUNNABLE};
\node[state,fill=c3!25]    (ru) at (8,2)       {RUNNING};
\node[state,fill=c4!25]    (wa) at (4,2)       {WAITING};
\node[state,fill=cred!25]  (zo) at (4,0)       {ZOMBIE};

\draw[arr] (un) -- node[above,font=\scriptsize]{create\_task} (cr);
\draw[arr] (cr) -- node[above,font=\scriptsize]{create\_process} (rn);
\draw[arr] (rn) -- node[right,font=\scriptsize]{scheduler} (ru);
\draw[arr] (ru) to[bend right=15] node[left,font=\scriptsize]{yield} (rn);
\draw[arr] (ru) -- node[above,font=\scriptsize]{suspend} (wa);
\draw[arr] (wa) -- node[right,font=\scriptsize]{wakeup} (rn);
\draw[arr] (ru) -- node[right,font=\scriptsize]{terminate} (zo);
\draw[arr,bend right=25] (zo) to node[above,font=\scriptsize]{free\_resources} (un);
\end{tikzpicture}
\end{center}

(EDOS define también \texttt{TERMINATED} en \texttt{task.h:15}, pero
en la práctica no la usa el flujo principal: las tareas pasan
directo de \texttt{RUNNING} a \texttt{ZOMBIE}.)

\subsection*{Walkthrough: \texttt{scheduler} (\texttt{task.c:90-118})}

\begin{walkthrough}
\begin{lstlisting}[style=cstyle]
void scheduler(void) {
    int cpu_id = cpuid();
    cpus_state[cpu_id].task = NULL;
    while (1) {
        enable_interrupts();
        for (int i=0; i<TASK_MAX; i++) {
            acquire(&tasks[i].lock);
            if (tasks[i].state == ZOMBIE) {
                free_resources(&tasks[i]);     // libera kstack y marca UNUSED
            } else if (tasks[i].state == RUNNABLE) {
                struct task* next = &tasks[i];
                next->state  = RUNNING;
                next->cpu_id = cpu_id;
                next->ticks  = QUANTUM;
                cpus_state[cpu_id].task = next;

                context_switch(&cpus_state[cpu_id].ctx, &next->ctx);

                // cuando next cede, volvemos acá
                next->cpu_id = -1;
                cpus_state[cpu_id].task = NULL;
            }
            release(&tasks[i].lock);
        }
    }
}
\end{lstlisting}

\textbf{Observaciones:}

\begin{itemize}[leftmargin=*,nosep]
    \item Cada CPU tiene su propio ``thread de scheduler'' con
        contexto en \texttt{cpus\_state[cpu].ctx}.
    \item El recorrido es \textbf{round-robin sobre el arreglo
        \texttt{tasks[]}}, no por prioridad ni por cola de listos
        separada.
    \item Al hacer \texttt{context\_switch}, el ``ra'' guardado del
        scheduler le permite volver a la línea siguiente cuando la
        tarea ceda.
    \item Mientras tiene la CPU una tarea, su \texttt{lock} está
        tomado (\textbf{importante}: \texttt{yield} libera tras el
        \texttt{sched}, ver abajo).
\end{itemize}
\end{walkthrough}

\subsection*{\texttt{yield}, \texttt{sched} y la pareja con \texttt{scheduler}}

\begin{lstlisting}[style=cstyle]
/* task.c:140-158 */
void yield(void) {
    struct task *t = cpus_state[cpuid()].task;
    acquire(&t->lock);
    t->state = RUNNABLE;
    sched(t);             // salta al scheduler, manteniendo el lock
    release(&t->lock);    // al volver, libera el lock
}

/* task.c:126-137 */
static void sched(struct task* current) {
    int cpu_id = cpuid();
    int irq_status = cpus_state[cpu_id].irq_enabled;
    context_switch(&current->ctx, &cpus_state[cpu_id].ctx);
    cpus_state[cpuid()].irq_enabled = irq_status;
}
\end{lstlisting}

\textbf{La danza del lock}:

\begin{enumerate}[nosep,leftmargin=*]
    \item \texttt{yield} entra y \emph{toma} el lock de la tarea.
    \item Cambia \texttt{state} a \texttt{RUNNABLE}.
    \item Llama a \texttt{sched}, que hace \texttt{context\_switch}
        al scheduler.
    \item El scheduler ``retorna'' en el medio de \texttt{for},
        donde luego va a \texttt{release(\&tasks[i].lock)}: ese es
        \emph{nuestro} lock.
    \item Más tarde, otro \texttt{context\_switch} nos trae de
        vuelta dentro de \texttt{sched} $\to$ \texttt{yield} $\to$
        \texttt{release(\&t->lock)}: ahora \emph{el scheduler nos
        liberó} y nosotros liberamos también. Cuidado: la cuenta
        de \texttt{noff} (push\_irq\_off) debe alinearse.
\end{enumerate}

\subsection*{Walkthrough: \texttt{create\_task}}

\begin{walkthrough}
\textbf{\texttt{task.c:53-79}}

\begin{lstlisting}[style=cstyle]
struct task* create_task(char* name, task_function f) {
    struct task *task = NULL;
    for (int i = 0; i < TASK_MAX && !task; i++) {
        acquire(&tasks[i].lock);
        if (tasks[i].state == UNUSED) {
            task = &tasks[i];
            memset(&task->ctx, 0, sizeof(struct context));
            task->ctx.ra = (address) f;                      // (a)
            task->kstack = alloc_page();                     // (b)
            task->ctx.sp = (address) (task->kstack + PAGE_SIZE);
            task->tid    = ++last_tid;
            task->pid    = 0;
            task->pgtbl  = kernel_pgtbl;                     // (c)
            task->state  = CREATED;
            strcpy(task->name, name);
        }
        release(&tasks[i].lock);
    }
    return task;
}
\end{lstlisting}

\textbf{(a)} El truco está acá: configuramos el ``return address''
del contexto guardado para que apunte a \texttt{f}. Cuando el
scheduler haga \texttt{context\_switch(\&scheduler\_ctx,
\&task->ctx)}, lo último que carga es \texttt{ra = f}, e
inmediatamente hace \texttt{ret} (la instrucción de retorno de
\texttt{context\_switch} en \texttt{arch.s:122}). Como \texttt{ret}
en RISC-V es \texttt{jr ra}, terminamos saltando a \texttt{f} ``como
si \texttt{f} hubiera sido llamada''.

\textbf{(b)} Reservamos una página para kernel stack y dejamos
\texttt{sp} en el tope (RISC-V crece para abajo).

\textbf{(c)} Tarea de kernel: comparte el \texttt{kernel\_pgtbl}.
Un proceso va a tener su propio pgtbl, asignado más tarde por
\texttt{exec}.
\end{walkthrough}

\subsection*{\texttt{suspend}, \texttt{wakeup}, \texttt{sleep}}

\begin{walkthrough}
\textbf{\texttt{task.c:163-183}}

\begin{lstlisting}[style=cstyle]
void suspend(void* condition, spinlock* lk) {
    struct task* task = current_task();
    acquire(&task->lock);
    release(lk);                  // liberamos el lock del recurso
    task->wait_condition = condition;
    task->state = WAITING;
    sched(task);                  // ceder
    task->wait_condition = 0;
    release(&task->lock);
    acquire(lk);                  // reasumimos el lock del recurso
}
\end{lstlisting}

Es el patrón clásico (xv6, Linux): \emph{atomically release+sleep}.
El truco para que el \texttt{wakeup} no se pierda es que
\texttt{lk} se libera \emph{con \texttt{task->lock} tomado}; el
\texttt{wakeup}, cuando recorre la tabla, intenta tomar
\texttt{task->lock}, así que no puede despertarnos antes de que
ya estemos en \texttt{WAITING} dentro de \texttt{sched}.

\texttt{wakeup} (\texttt{task.c:186-200}) recorre todas las tareas,
toma su \texttt{lock} y cambia a \texttt{RUNNABLE} las que coincidan
en \texttt{wait\_condition}.

\texttt{sleep(n)} (\texttt{task.c:203-212}) es un caso particular:
deja en \texttt{sleep\_ticks} el contador y se suspende con
condición \texttt{\&ticks}. El \texttt{inc\_ticks} en cada timer
decrementa \texttt{sleep\_ticks} y, al llegar a 0, pasa a
\texttt{RUNNABLE}.
\end{walkthrough}

\subsection*{Errores típicos en \texttt{task.c}}

\begin{itemize}[leftmargin=*,nosep]
    \item Olvidar el \texttt{release(\&t->lock)} en algún path:
        deja el lock tomado para siempre.
    \item Hacer \texttt{wakeup(t)} sin tener \texttt{t->lock} (en
        EDOS lo toma adentro). Si lo hicieras desde fuera con el
        lock ya tomado, deadlock.
    \item Cambiar el \texttt{state} fuera del lock.
    \item Asumir que después de \texttt{sched(t)} estamos en la
        misma CPU. Podemos haber migrado.
\end{itemize}

% ---------------------------------------------------------------------
\section{Procesos de usuario (task.c + vm.c + arch.s)}

\subsection*{Espacio de direcciones del proceso}

\texttt{PROC\_MIN\_VA = 0}, \texttt{PROC\_MAX\_VA = 0x80000000}
(\texttt{task.h:19-20}). El stack va al final, una sola página:

\begin{center}
\begin{tikzpicture}[
    every node/.style={font=\ttfamily\small},
    sec/.style={draw,minimum width=4.5cm,minimum height=0.7cm,align=center}
]
\node[sec,fill=c1!25] (text)  at (0,5)        {código + datos\\(desde \texttt{0x0})};
\node[sec,fill=cfree!20,below=0pt of text,minimum height=2.2cm] (gap) {free space \\ (sin mapear)};
\node[sec,fill=c3!25,below=0pt of gap]        (stk)  {stack \\ (1 página)};
\node[anchor=west,xshift=0.5cm] at (text.east) {\texttt{0x00000000}};
\node[anchor=west,xshift=0.5cm] at (gap.east)  {};
\node[anchor=west,xshift=0.5cm] at (stk.east)  {\texttt{PROC\_MAX\_VA -- 4\,KiB}};
\node[anchor=west,xshift=0.5cm] at (stk.south east) {\texttt{0x80000000} (excluido)};
\end{tikzpicture}
\end{center}

\subsection*{Creación del primer proceso}

\begin{walkthrough}
\textbf{\texttt{task.c:424-441} (\texttt{create\_process})}

\begin{lstlisting}[style=cstyle]
bool create_process(char *path) {
    struct task* t = create_task(path, start_process); // (a)
    if (!t) panic("Create process: No tasks free slot!");

    t->pid = ++last_pid;                                // (b)
    if (!exec(t, path, 0)) {                            // (c)
        release(&t->lock);
        return false;
    }
    t->state = RUNNABLE;                                // (d)
    return true;
}
\end{lstlisting}

\textbf{(a)} Crear un \texttt{task} con \texttt{start\_process} como
``ra''. Cuando el scheduler lo elija por primera vez,
\texttt{context\_switch} le entregará el control a
\texttt{start\_process}.

\textbf{(b)} Asignar PID (\texttt{++last\_pid}).

\textbf{(c)} \texttt{exec} carga el programa.

\textbf{(d)} \texttt{RUNNABLE}: ya es elegible.
\end{walkthrough}

\subsection*{Walkthrough: \texttt{exec}}

\begin{walkthrough}
\textbf{\texttt{task.c:376-421}}

\begin{lstlisting}[style=cstyle]
bool exec(struct task* task, char* path, char* args[]) {
    struct trap_frame *tf = task_trap_frame_address(task);

    pte* new_pgtbl = (pte*) alloc_page();
    if (!new_pgtbl) goto error;

    memcpy(new_pgtbl, kernel_pgtbl, PAGE_SIZE);          // (1)
    if (!load_program(new_pgtbl, path)) goto error;       // (2)

    uint8 *ustack = (uint8 *) alloc_page();
    if (!ustack) goto error;
    map_page(new_pgtbl, PROC_MAX_VA - PAGE_SIZE, (paddr) ustack,
             PAGE_R | PAGE_W | PAGE_U);                   // (3)

    vaddr sp = setup_main_args(task, ustack + PAGE_SIZE, args); // (4)

    tf->sp = sp;                                          // (5)
    tf->pc = 0;                                           // (6)

    if (is_process(task)) {                               // (7)
        unmap_region(task->pgtbl, 0, PROC_MAX_VA, true);
        free_page(task->pgtbl);
    }
    task->pgtbl  = new_pgtbl;
    task->ctx.sp = (size_t) tf;                           // (8)
    return sp;
error:
    free_page(new_pgtbl);
    return false;
}
\end{lstlisting}

\textbf{(1)} \texttt{memcpy(new\_pgtbl, kernel\_pgtbl, PAGE\_SIZE)}.
Copia el directorio raíz del kernel en el del proceso. \emph{Las
entradas de hoja del kernel (sus PTs internas) son compartidas con
toda otra page table de proceso}. Por eso un page fault en código
de \texttt{s\_trap} jamás ocurre: está mapeado en cualquier proceso.

\textbf{(2)} \texttt{load\_program} (\texttt{task.c:291-313}) hace
\texttt{efs\_file(path)}, asigna páginas para todo el contenido y
las mapea con \texttt{R|W|X|U} a partir de \texttt{VA=0}.

\textbf{(3)} Mapea \emph{una sola página de stack} con permisos
\texttt{R|W|U} en \texttt{PROC\_MAX\_VA - 4\,KiB}. Sin \texttt{U}
el proceso no podría ni escribir su SP.

\textbf{(4)} \texttt{setup\_main\_args} (\texttt{task.c:335-363})
copia los strings de \texttt{args[]} al stack y arma el array
\texttt{argv[]}. Llena \texttt{tf->a0 = argc} y
\texttt{tf->a1 = \&argv}.

\textbf{(5)(6)} Punto de retorno a usuario: \texttt{pc=0} (entry
point del programa) y \texttt{sp = top del stack}.

\textbf{(7)} Si el task ya era un proceso (caso de un \texttt{exec}
sobre un proceso vivo), desmapea su anterior pgtbl y libera.

\textbf{(8)} \texttt{task->ctx.sp = (size\_t) tf}. Esto es elegante:
deja el kernel-mode SP del task apuntando al trap frame. Cuando el
scheduler haga \texttt{context\_switch} a este task y caigamos en
\texttt{start\_process}, \texttt{sp} ya estará sobre el trap frame
listo para que \texttt{return\_to\_user\_mode} lea
\texttt{task->kstack + PAGE\_SIZE - sizeof(tf)}.
\end{walkthrough}

\subsection*{\texttt{start\_process} y la primera vuelta a U-mode}

\begin{walkthrough}
\textbf{\texttt{task.c:366-372}}

\begin{lstlisting}[style=cstyle]
static void start_process(void) {
    struct task *task = current_task();
    release(&task->lock);          // soltamos el lock que tomó el scheduler
    return_to_user_mode();          // configurar CSRs y u_trap_ret
}
\end{lstlisting}

Recordá: el \emph{scheduler} cuando hace \texttt{context\_switch} a
un task \texttt{RUNNABLE} ya tiene \texttt{task->lock} tomado. La
``promesa'' es que la tarea, al recibir el control, va a soltarlo.

Después de \texttt{return\_to\_user\_mode}, el flujo termina en
\texttt{u\_trap\_ret} (\texttt{trap.s:150-192}) que pone \texttt{sp}
al trap frame, restaura los registros y hace \texttt{sret}. El
\texttt{sret} salta a \texttt{tf->pc = 0}, que es la entry
\texttt{start} del binario de usuario (definida primero en
\texttt{user.ld} con \texttt{ENTRY(start)}).
\end{walkthrough}

\subsection*{Diagrama: ciclo de vida de un proceso}

\begin{center}
\begin{tikzpicture}[
    every node/.style={font=\ttfamily\small,align=center},
    box/.style={draw,rounded corners,minimum width=3.2cm,minimum height=0.8cm},
    arr/.style={->,thick}
]
\node[box,fill=c1!25] (cm) at (0, 5) {kernel\_main \\ \scriptsize (hart 0)};
\node[box,fill=c2!25] (cp) at (0, 3) {create\_process};
\node[box,fill=c3!25] (sc) at (0, 1) {scheduler \\ \scriptsize (cualquier hart)};
\node[box,fill=c4!25] (sp) at (6, 1) {start\_process};
\node[box,fill=c4!25] (rt) at (6, 3) {return\_to\_user\_mode};
\node[box,fill=cred!25] (us) at (6, 5) {init.c::start \\ \scriptsize (U-mode)};
\draw[arr] (cm) -- node[right,font=\scriptsize]{create\_process("init")} (cp);
\draw[arr] (cp) -- node[right,font=\scriptsize]{exec + state=RUNNABLE} (sc);
\draw[arr] (sc) -- node[above,font=\scriptsize]{context\_switch} (sp);
\draw[arr] (sp) -- (rt);
\draw[arr] (rt) -- node[right,font=\scriptsize]{u\_trap\_ret + sret} (us);
\end{tikzpicture}
\end{center}

% ---------------------------------------------------------------------
\section{Syscalls (syscall.c, syscall.h)}

\subsection*{Convención}

Los binarios de usuario implementan cada syscall en
\texttt{user/usys.s} cargando el número en \texttt{a7} y haciendo
\texttt{ecall}:

\begin{lstlisting}[style=asmstyle]
.global getpid
getpid:
    li a7, 1     # SYS_GETPID
    ecall
    ret
\end{lstlisting}

Los argumentos viajan en \texttt{a0--a6}. El resultado vuelve en
\texttt{a0}.

\subsection*{Tabla de syscalls}

\begin{center}
\small
\begin{tabular}{@{}clp{8cm}@{}}
\toprule
\# & Nombre & Comportamiento \\
\midrule
0 & \texttt{exit(int code)}      & Marca \texttt{killed=true}; el trap handler lo cosechará en \texttt{terminate}. \\
1 & \texttt{getpid()}            & Devuelve \texttt{task->pid}. \\
2 & \texttt{console\_puts(s)}    & Traduce la VA del string a kernel y hace \texttt{printf("\%s", \dots)}. \\
3 & \texttt{console\_putc(c)}    & \texttt{console\_putc(arg0)}. \\
4 & \texttt{console\_getc()}     & Lee 1 carácter (no bloqueante; devuelve 0 si no hay). \\
5 & \texttt{sleep(n)}            & Suspende la tarea por \texttt{n} ticks. \\
\bottomrule
\end{tabular}
\end{center}

\subsection*{Walkthrough: el dispatcher (\texttt{syscall.c:73-84})}

\begin{walkthrough}
\begin{lstlisting}[style=cstyle]
void syscall(struct task *task) {
    struct trap_frame *tf = task_trap_frame_address(task);
    uint n = syscall_number(tf);              // tf->a7
    if (n < SYSCALLS && syscalls_table[n]) {
        int result = syscalls_table[n](task);
        syscall_put_result(tf, result);       // tf->a0 = result
    } else {
        syscall_put_result(tf, -1);
    }
}
\end{lstlisting}

Notar que el \texttt{syscall} \emph{escribe directamente en el trap
frame}. Cuando \texttt{user\_trap} retorna y \texttt{u\_trap\_ret}
restaura registros, \texttt{a0} contiene el resultado.

Después de retornar, \texttt{user\_trap} llama a
\texttt{skip\_trap\_instruction(tf)} (\texttt{arch.h:333}) que hace
\texttt{tf->pc += 4}. Sin esto, el \texttt{sret} volvería al mismo
\texttt{ecall} y caeríamos en un loop infinito.
\end{walkthrough}

\subsection*{Validación de punteros}

\textbf{Importante}: \texttt{sys\_console\_puts}
(\texttt{syscall.c:22-29}) traduce la VA del usuario a una dirección
kernel:

\begin{lstlisting}[style=cstyle]
int sys_console_puts(struct task *task) {
    struct trap_frame *tf = task_trap_frame_address(task);
    size_t str  = syscall_arg(tf, 0);
    char *kaddr = (char *) va2kernel_address(task->pgtbl, (vaddr)str);
    printf("%s", kaddr);
    return 0;
}
\end{lstlisting}

\begin{importante}
La validación en EDOS \textbf{es mínima}: si \texttt{kaddr} resulta
0 (página no mapeada en el proceso), el \texttt{printf("\%s", 0)}
va a crashear. Tampoco se chequea longitud, alineamiento ni que
toda la cadena viva en una región accesible. Esto es deliberado:
EDOS muestra el \emph{esqueleto}, no la validación industrial. En
xv6 lo hace \texttt{argstr/argint/fetchaddr}.
\end{importante}

% ---------------------------------------------------------------------
\section{Sistema de archivos EFS (efs.c, efs.h, efsfiles.c)}

\subsection*{Qué provee}

Un sistema de archivos \textbf{en RAM}, \textbf{embebido en el
binario del kernel}, \textbf{de sólo lectura}. Cada archivo es un
arreglo de bytes con un nombre y un tipo (\texttt{EFS\_FILE\_PROGRAM}
o \texttt{EFS\_FILE\_DATA}).

\subsection*{Tabla de archivos}

\begin{lstlisting}[style=cstyle]
struct file {
    char            *name;
    unsigned char   type;        // PROGRAM o DATA
    unsigned char   *data;
    unsigned int    length;
};
extern struct file efs_files_table[];
\end{lstlisting}

\texttt{efsfiles.c} (\textbf{generado} por \texttt{user/mkefs.sh}
con \texttt{xxd -i}) contiene:

\begin{itemize}[nosep,leftmargin=*]
    \item \texttt{unsigned char init\_bin[]} — bytes binarios del
        ejecutable \texttt{init}.
    \item \texttt{unsigned char README[]} — contenido del README.
    \item \texttt{efs\_files\_table[]} = \{(\texttt{init},
        \texttt{PROGRAM}, \texttt{init\_bin}, len), (\texttt{README},
        \texttt{DATA}, \texttt{README}, len), \{0\}\}.
\end{itemize}

\subsection*{API}

Sólo dos cosas, y una de ellas con una línea:

\begin{lstlisting}[style=cstyle]
struct file* efs_file(char *file_name) {
    for (int i=0; efs_files_table[i].name; i++)
        if (strcmp(efs_files_table[i].name, file_name) == 0)
            return efs_files_table + i;
    return 0;
}
\end{lstlisting}

\subsection*{Cómo \texttt{exec} lo usa}

En \texttt{task.c:291-313} (\texttt{load\_program}):

\begin{enumerate}[nosep,leftmargin=*]
    \item \texttt{efs\_file(path)} devuelve el \texttt{file*} de
        \texttt{init} (o falla si el path no existe / no es PROGRAM).
    \item Por cada bloque de hasta 4\,KiB, hace
        \texttt{alloc\_page}, copia con \texttt{memcpy} y
        \texttt{map\_page} a la VA correspondiente.
\end{enumerate}

\subsection*{Diferencias con xv6fs}

EFS \textbf{no tiene}:

\begin{itemize}[nosep,leftmargin=*]
    \item Inodes ni superblock.
    \item Bloques en disco (todo en RAM).
    \item Directorios (es una lista plana).
    \item Permisos.
    \item Escritura, append, truncate, delete.
\end{itemize}

Esto es a propósito: la consigna no es ``filesystem completo''
sino ``manera de tener un binario de usuario disponible al boot''.
xv6fs (que sí tiene todo lo anterior) se cubre en talleres
posteriores en otras versiones del curso.

% ---------------------------------------------------------------------
\section{Programas de usuario (user/)}

\subsection*{Qué hay}

\begin{itemize}[nosep,leftmargin=*]
    \item \texttt{init.c} — el ``programa init'' que es el primer
        (y único, en 08-process) proceso de usuario.
    \item \texttt{edoslib.c}/\texttt{edoslib.h} — la mini-libc de
        EDOS. Provee \texttt{printf}, \texttt{strcpy}, \texttt{memset}, etc.,
        más la entrada \texttt{start()}.
    \item \texttt{usys.s} — un stub por syscall que carga
        \texttt{a7} y hace \texttt{ecall}.
    \item \texttt{user.ld} — linker script con \texttt{ENTRY(start)}
        y \texttt{. = 0;}.
    \item \texttt{Makefile} — compila a rv32 estático y empaqueta.
    \item \texttt{mkefs.sh} — convierte los binarios en
        \texttt{efsfiles.c}.
    \item \texttt{README} — archivo de datos para probar EFS.
\end{itemize}

\subsection*{Cómo se compila el ejecutable}

\begin{lstlisting}[style=shellstyle]
riscv64-unknown-elf-gcc -O -Wall -mcmodel=medany -march=rv32ima -mabi=ilp32 \
  -ffreestanding -nostdlib -fno-common -fno-omit-frame-pointer \
  -fno-stack-protector -nostartfiles -static \
  -Wl,-Tuser.ld,--build-id=none -o init -e start \
  edoslib.o usys.o init.c
riscv64-unknown-elf-objcopy -S -O binary init init.bin
\end{lstlisting}

\textbf{Notas:}

\begin{itemize}[leftmargin=*,nosep]
    \item \texttt{rv32ima} = rv32 base + multiplicación + atómicos.
    \item \texttt{nostartfiles} suprime el \texttt{crt0.o} de
        GCC: usamos \emph{nuestra} \texttt{start} en \texttt{edoslib.c:14}.
    \item \texttt{-e start} explicita el entry point.
    \item \texttt{objcopy -O binary} elimina la cabecera ELF; el
        binario que termina en EFS son \emph{instrucciones puras}
        listas para mapear en VA=0.
\end{itemize}

\subsection*{\texttt{start} y el ciclo de vida}

\begin{lstlisting}[style=cstyle]
extern int main();
void start(void) { exit(main()); }
\end{lstlisting}

(\texttt{edoslib.c:14-17}.) La primera función definida en el
\texttt{.text}, garantizada por \texttt{KEEP(*(.text.start))} en
\texttt{user.ld}. Aunque acá usamos sólo la convención del nombre,
en proyectos más grandes habría una sección \texttt{.text.start}
explícita.

% =====================================================================
\part{Diferencias y similitudes con xv6}
% =====================================================================

\section{Tabla comparativa}

\begin{center}
\footnotesize
\begin{longtable}{@{}p{3.5cm}p{6cm}p{6cm}@{}}
\toprule
Aspecto & \textbf{EDOS (08-process)} & \textbf{xv6-riscv} \\
\midrule
ISA target           & rv32ima (\texttt{ilp32}) & rv64 (\texttt{lp64}) \\
Paginación           & Sv32, 2 niveles, página 4\,KiB & Sv39, 3 niveles, página 4\,KiB \\
Flags PTE expuestos  & V, R, W, X, U & V, R, W, X, U, A, D, G \\
Layout kernel        & 1:1 desde \texttt{0x80000000}; MMIO mapeado 1:1; sin trampoline & 1:1 + trampoline en la VA más alta (en cada proceso) \\
Trampoline           & No usa página separada; \texttt{s\_trap}/\texttt{u\_trap} viven en \texttt{.text} y se ven gracias a la copia de \texttt{kernel\_pgtbl} & Página dedicada \texttt{TRAMPOLINE} mapeada en cada proceso y en el kernel \\
Trap frame           & En la cima del kernel-stack del task & Página \texttt{trapframe} per-proc mapeada vía VA fija \\
PCB                  & \texttt{struct task} de \texttt{task.c}; UNIFICA threads de kernel y procesos & \texttt{struct proc}; tienen \texttt{trapframe} y \texttt{context} separados \\
Estados              & UNUSED, CREATED, RUNNABLE, RUNNING, WAITING, ZOMBIE & UNUSED, USED, SLEEPING, RUNNABLE, RUNNING, ZOMBIE \\
Scheduler            & Round-robin sobre \texttt{tasks[]}, quantum=2 ticks, sin prioridades & Round-robin con sleep-locks; idéntica filosofía \\
Context switch       & \texttt{context\_switch} en \texttt{arch.s}, 14 callee-saved (rv32) & \texttt{swtch.S}, 14 callee-saved (rv64) \\
Locking              & \texttt{spinlock} (\texttt{unsigned int}) con \texttt{push/pop\_irq\_off} & \texttt{spinlock} (con holder, name) + \texttt{sleeplock} encima \\
Sistema de archivos  & EFS (embebido, RAM, read-only, sin directorios) & xv6fs (en disco, journaled, inodes, directorios) \\
Syscalls             & 6 (\texttt{exit, getpid, console\_puts, console\_putc, console\_getc, sleep}) & $\sim$22 (incluye \texttt{fork, exec, wait, read, write, open, close, pipe, fstat, dup, sbrk, kill, mkdir, link, unlink, chdir}) \\
Convención syscall   & \texttt{a7} = número, \texttt{a0--a6} = args, \texttt{a0} = ret (idéntica) & Idéntica \\
Driver UART          & NS16550A, polling TX, RX configurable & NS16550A, IRQ-driven RX, ring buffer \\
SMP                  & Soporta hasta NCPU=4, pero \texttt{Makefile} arranca \texttt{-smp 1} & Plenamente SMP; un hart 0 inicia y los demás esperan \\
M-mode               & EDOS hace el rol de OpenSBI: setea PMP, deleg, MPP, timer & xv6 con \texttt{-bios default} se apoya en OpenSBI; el código RISC-V en xv6 entra ya en S-mode \\
\bottomrule
\end{longtable}
\end{center}

\section{Discusión: el porqué pedagógico de cada decisión}

\paragraph{rv32 vs.\ rv64.}
xv6-riscv eligió rv64 porque emula servidores modernos. EDOS optó
por rv32 para que (a) las direcciones quepan en \texttt{uint32}, (b)
las traducciones de páginas tengan sólo \textbf{dos niveles}, (c)
los listados de ensamblador usen \texttt{lw}/\texttt{sw} (más
intuitivos) y (d) la consigna ``imprimir un puntero'' no asuste con
8 dígitos hexa.

\paragraph{Sv32 vs.\ Sv39.}
Mismo motivo: dos niveles permiten dibujar la traducción completa
en un slide. Sv39 tiene 3 niveles + páginas de 1\,GiB y 2\,MiB que
quitan foco.

\paragraph{Sin trampoline page.}
xv6 mapea una página de código (\texttt{trampoline.S}) en la VA más
alta tanto en kernel como en cada proceso. Esa página tiene
\texttt{uservec}/\texttt{userret}. EDOS toma un atajo: la copia del
directorio raíz del kernel en cada \texttt{pgtbl} de proceso
incluye las PTEs que apuntan a las páginas del kernel.text, y por
tanto \texttt{s\_trap}/\texttt{u\_trap} \emph{son} accesibles desde
el proceso (aunque sin la flag \texttt{U}, claro: sólo el hardware
en modo S puede saltar ahí).

\paragraph{Trap frame en el kernel-stack.}
xv6 reserva una página \texttt{trapframe} por proceso, en una VA
fija. En EDOS, el trap frame se construye \emph{empujando} 31
registros sobre el kernel-stack del task; la macro
\texttt{task\_trap\_frame\_address(t)} es simplemente
\texttt{t->kstack + PAGE\_SIZE - sizeof(struct trap\_frame)}.
Funciona porque la página del kernel-stack ya está accesible al
kernel desde cualquier proceso. La ventaja didáctica: un solo dato
para entender (``la cima del stack''), sin VA mágica adicional.

\paragraph{EFS vs.\ xv6fs.}
xv6 tiene un FS de disco completo (inodes, journal, directorios,
crash-safety). Para entender ``cómo se carga el init'' no hace
falta nada de eso: EFS pega los binarios al kernel con \texttt{xxd}.
A cambio el alumno aprende que ``filesystem'' no significa
``ext4''; significa ``mapeo nombre $\to$ bytes''.

\paragraph{Round-robin puro.}
EDOS no tiene prioridades ni MLFQ. La cátedra muestra esa
complejidad en los prácticos teóricos; el kernel se queda con la
forma más simple para que la atención esté en \texttt{context\_switch}
y la danza del lock con el scheduler.

% =====================================================================
\part{Banco de preguntas de defensa}
% =====================================================================

A continuación, 35 preguntas razonables agrupadas por tema, con
respuestas completas que citan archivo y línea para que las puedas
ir leyendo en paralelo.

\section{Conceptuales generales}

\subsection*{P1. ¿Qué es un trap y por qué nuestro kernel \emph{necesita} uno?}

\begin{respuesta}
Un \emph{trap} es cualquier transferencia involuntaria del flujo
de control desde el código actual hacia un handler en modo
supervisor. En RISC-V hay tres familias: \textbf{interrupciones}
asíncronas (timer, dispositivos), \textbf{excepciones} síncronas
(page fault, instrucción ilegal, división por cero) y la
\textbf{trap explícita} (\texttt{ecall}, que un proceso usa para
pedirle al kernel un servicio).

Sin traps, el kernel \textbf{no podría existir tal como lo
conocemos}: ningún proceso podría pedirle servicios al SO sin
ejecutar él mismo en modo privilegiado, y el SO no podría recuperar
la CPU de un proceso que no la cede voluntariamente. La existencia
de traps + dos modos (U/S) es la base de la separación entre
``programa de usuario'' y ``kernel''.

En EDOS los handlers de bajo nivel están en \texttt{trap.s}: tres
de ellos (\texttt{s\_trap}, \texttt{u\_trap}, \texttt{m\_trap}) y se
ramifican en C en \texttt{trap.c::kernel\_trap} y
\texttt{trap.c::user\_trap}.
\end{respuesta}

\subsection*{P2. ¿Qué es un trampoline y EDOS lo usa?}

\begin{respuesta}
En xv6 el \emph{trampoline} es una página de código mapeada en una
VA fija tanto en el espacio del kernel como en el de cada proceso.
Contiene \texttt{uservec} y \texttt{userret} (los pasajes
\texttt{U $\to$ S} y \texttt{S $\to$ U}). El nombre viene de que
``ambos espacios pueden saltar en él''.

\textbf{EDOS no usa una página trampoline separada.} En su lugar:

\begin{itemize}[leftmargin=*,nosep]
    \item Las rutinas \texttt{s\_trap} y \texttt{u\_trap} viven en
        el segmento \texttt{.text} del kernel
        (\texttt{trap.s:10,89}).
    \item \texttt{exec} copia el directorio raíz del kernel en cada
        \texttt{pgtbl} de proceso (\texttt{task.c:386}). Esa copia
        incluye las PTEs intermedias que llevan a la \texttt{.text}
        del kernel.
    \item Por lo tanto la VA donde está \texttt{u\_trap} es
        \emph{traducible} (a la misma PA, vía 1:1) desde cualquier
        proceso. Cuando ocurre el trap y el hardware salta a
        \texttt{stvec = u\_trap}, no hay un cambio de \texttt{satp}:
        el código del kernel ya está accesible.
\end{itemize}

La consecuencia es que \texttt{u\_trap} \textbf{no} cambia el
\texttt{satp} (es decir, sigue con el \texttt{pgtbl} del proceso)
mientras hace el trabajo bajo nivel; sólo cuando entra a
\texttt{user\_trap} se hace \texttt{set\_page\_table(kernel\_pgtbl)}
(\texttt{trap.c:126}).
\end{respuesta}

\subsection*{P3. ¿Qué hereda un task que viene del scheduler? ¿Cómo se entera el scheduler de dónde retomar?}

\begin{respuesta}
La pieza es \texttt{context\_switch} (\texttt{arch.s:90-122}). Se
invoca como \texttt{context\_switch(\&old, \&new)} con dos punteros
a \texttt{struct context}. La rutina:

\begin{enumerate}[leftmargin=*,nosep]
    \item Guarda en \texttt{*old} los registros callee-saved
        actuales: \texttt{ra, sp, s0--s11}.
    \item Restaura desde \texttt{*new} esos mismos registros.
    \item Hace \texttt{ret} (jr ra), que salta al \texttt{ra}
        recién restaurado.
\end{enumerate}

Por eso, cuando el scheduler hace \texttt{context\_switch(\&own,
\&task->ctx)}:

\begin{itemize}[leftmargin=*,nosep]
    \item Guardó su propio \texttt{ra} (la instrucción siguiente a
        la llamada) en \texttt{own.ra}. \emph{Esa es su marca de
        retorno.}
    \item Cargó \texttt{task->ctx.ra} y saltó allí. Si la tarea
        nunca corrió, \texttt{ra} apunta a \texttt{start\_process}
        (preparado en \texttt{create\_task}). Si ya corrió y luego
        cedió, \texttt{ra} apunta justo después del
        \texttt{context\_switch} \emph{de \texttt{sched}} que la
        cedió.
\end{itemize}

Lo importante: el control retorna al scheduler exactamente cuando
otro task hace \texttt{context\_switch(\&task->ctx,
\&cpu->ctx)}, y entonces el flujo retorna a la línea siguiente
\texttt{cpus\_state[cpu\_id].task = NULL;} en
\texttt{task.c:112-113}.
\end{respuesta}

\subsection*{P4. ¿Qué pasa con los stacks durante el cambio de contexto?}

\begin{respuesta}
\texttt{context\_switch} salva y restaura el \texttt{sp} como
cualquier otro callee-saved (\texttt{arch.s:93,109}). Entonces:

\begin{itemize}[leftmargin=*,nosep]
    \item Cuando se entra a \texttt{context\_switch}, estamos en el
        stack del que cede.
    \item Después del \texttt{lw sp, 4(a1)} estamos en el stack del
        próximo. El \texttt{ret} usa ese stack.
\end{itemize}

Es seguro porque ambas estructuras \texttt{context} fueron
preparadas con \texttt{sp = base del kstack + 4\,KiB} en algún
momento previo, y la zona de stack está siempre \emph{válidamente
mapeada} en el \texttt{kernel\_pgtbl} (todos los stacks viven en la
RAM 1:1).
\end{respuesta}

\subsection*{P5. ¿Cómo asegura el spinlock que no haya deadlock entre un trap y la sección crítica?}

\begin{respuesta}
Porque \texttt{acquire} llama a \texttt{push\_irq\_off}
(\texttt{spinlock.c:7-17}) antes de hacer el test-and-set. Esto
deshabilita las IRQs en la CPU local. Mientras tenemos el lock, no
puede llegar un trap del timer que invoque \texttt{yield} y nos
deje girando para siempre. \texttt{pop\_irq\_off} sólo rehabilita
cuando el último \texttt{release} (anidamiento contado en
\texttt{noff}) lo permita.

En SMP, IRQs deshabilitadas en una CPU no garantizan exclusión
entre CPUs (eso sí lo hace el spinlock); pero sí garantizan que la
propia CPU no se interrumpa a sí misma estando dentro de la
sección crítica.
\end{respuesta}

\section{Caminá el código de…}

\subsection*{P6. Caminá el \texttt{exec} desde \texttt{create\_process} hasta el primer \texttt{sret}.}

\begin{respuesta}
\textbf{Paso a paso}, con citas:

\begin{enumerate}[leftmargin=*]
    \item \texttt{create\_process("init")} (\texttt{task.c:424-441})
        llama a \texttt{create\_task("init", start\_process)}.
    \item \texttt{create\_task} (\texttt{task.c:53-79}) busca un
        slot \texttt{UNUSED} en \texttt{tasks[]}, le asigna kstack,
        \texttt{ctx.ra=start\_process}, \texttt{ctx.sp=kstack+4K},
        \texttt{pgtbl=kernel\_pgtbl}, estado \texttt{CREATED}.
    \item \texttt{create\_process} asigna \texttt{pid} y llama a
        \texttt{exec(t, "init", 0)}.
    \item \texttt{exec} (\texttt{task.c:376-421}):
        \begin{enumerate}[nosep,label=\alph*)]
            \item \texttt{alloc\_page} para el nuevo \texttt{pgtbl}.
            \item \texttt{memcpy(new\_pgtbl, kernel\_pgtbl,
                PAGE\_SIZE)} — copia las entradas raíz del kernel.
            \item \texttt{load\_program(new\_pgtbl, "init")}
                (\texttt{task.c:291-313}) llama
                \texttt{efs\_file("init")} y mapea cada bloque a
                VAs desde 0 con \texttt{R|W|X|U}.
            \item \texttt{alloc\_page} para el user stack y
                \texttt{map\_page} a \texttt{PROC\_MAX\_VA-4\,KiB}
                con \texttt{R|W|U}.
            \item \texttt{setup\_main\_args} (no hay args para init,
                pero igual configura \texttt{a0=0}, \texttt{a1=PROC\_MAX\_VA}).
            \item Setea en el trap frame: \texttt{tf->pc=0},
                \texttt{tf->sp=top stack user}.
            \item Asigna \texttt{task->pgtbl = new\_pgtbl} y
                \texttt{task->ctx.sp = (size\_t) tf}.
        \end{enumerate}
    \item \texttt{create\_process} pone \texttt{state=RUNNABLE}.
    \item En algún hart, el \texttt{scheduler} (\texttt{task.c:90-118})
        encuentra el task \texttt{RUNNABLE}, lo elige y hace
        \texttt{context\_switch} a \texttt{task->ctx}.
    \item Como \texttt{ctx.ra = start\_process}, el \texttt{ret} de
        \texttt{context\_switch} cae en \texttt{start\_process}
        (\texttt{task.c:366-372}) con \texttt{sp} apuntando al trap
        frame en la cima del kstack.
    \item \texttt{start\_process} libera el lock del task y llama
        \texttt{return\_to\_user\_mode()}.
    \item \texttt{return\_to\_user\_mode} (\texttt{trap.c:72-101}):
        \begin{enumerate}[nosep,label=\alph*)]
            \item \texttt{disable\_interrupts()}.
            \item \texttt{stvec = u\_trap}.
            \item \texttt{sscratch = task->kstack+4\,KiB}.
            \item \texttt{sstatus.SPP=0; SPIE=1}.
            \item \texttt{satp = task->pgtbl} (con
                \texttt{sfence.vma}).
            \item \texttt{u\_trap\_ret(\&trap\_frame)}.
        \end{enumerate}
    \item \texttt{u\_trap\_ret} (\texttt{trap.s:150-192}) pone
        \texttt{sp=trap\_frame}, restaura \texttt{sepc} desde
        \texttt{tf->pc=0}, restaura todos los registros (incluido
        \texttt{sp=tf->sp}) y hace \texttt{sret}.
    \item \texttt{sret} salta a la VA 0 \textbf{en U-mode}, donde
        está la primera instrucción de \texttt{start} (entry de
        \texttt{init} según \texttt{user.ld}).
\end{enumerate}
\end{respuesta}

\subsection*{P7. Caminá la \texttt{getpid()} desde el espacio de usuario.}

\begin{respuesta}
\begin{enumerate}[leftmargin=*]
    \item En \texttt{init.c}, \texttt{getpid()} es una declaración
        \texttt{extern int getpid(void)} de \texttt{edoslib.h:24}.
        Su definición está en \texttt{user/usys.s:13-17}:
\begin{lstlisting}[style=asmstyle]
getpid:
    li a7, 1       # SYS_GETPID
    ecall
    ret
\end{lstlisting}
    \item \texttt{ecall} en U-mode genera un trap con
        \texttt{scause=8} (\texttt{SYSCALL} en \texttt{arch.h:73}).
    \item El hardware salta a \texttt{stvec = u\_trap}
        (\texttt{trap.s:89}).
    \item \texttt{u\_trap} guarda registros en el kstack del
        proceso, llama \texttt{user\_trap()}.
    \item \texttt{user\_trap} (\texttt{trap.c:107-181}) en el
        \texttt{case SYSCALL} (linea 146):
        \begin{enumerate}[nosep,label=\alph*)]
            \item \texttt{enable\_interrupts()}.
            \item \texttt{syscall(task)} (en \texttt{syscall.c:73}).
            \item \texttt{skip\_trap\_instruction(tf)}.
        \end{enumerate}
    \item \texttt{syscall} lee \texttt{n = tf->a7 = 1}, indexa
        \texttt{syscalls\_table[1]} y llama \texttt{sys\_getpid(task)}
        (\texttt{syscall.c:16-19}), que devuelve \texttt{task->pid}.
    \item \texttt{syscall\_put\_result(tf, result)} pone
        \texttt{tf->a0 = task->pid}.
    \item \texttt{user\_trap} retorna; \texttt{return\_to\_user\_mode}
        y \texttt{u\_trap\_ret} restauran los registros del frame
        (\texttt{a0} ahora tiene el PID) y hacen \texttt{sret} a
        \texttt{tf->pc + 4} (\texttt{ecall} es 4 bytes).
    \item En U-mode, el \texttt{ret} de \texttt{usys.s} salta de
        vuelta al llamador, que ve el PID en \texttt{a0} según ABI
        \texttt{ilp32}.
\end{enumerate}
\end{respuesta}

\subsection*{P8. Caminá un \emph{page fault} del proceso.}

\begin{respuesta}
\begin{enumerate}[leftmargin=*]
    \item El proceso accede a una VA no mapeada (o con flags
        incorrectos: ej.\ escribir una página R+X).
    \item La MMU genera la excepción. \texttt{scause} toma uno de:
        \texttt{INSTRUCTION\_PAGE\_FAULT (0xc)},
        \texttt{LOAD\_PAGE\_FAULT (0xd)} o
        \texttt{STORE\_PAGE\_FAULT (0xf)} (\texttt{arch.h:75-79}).
        \texttt{stval} contiene la VA que falló y \texttt{sepc} la
        instrucción que la disparó.
    \item Salta a \texttt{stvec = u\_trap}.
    \item \texttt{u\_trap} guarda el trap frame y llama
        \texttt{user\_trap}.
    \item \texttt{user\_trap} entra a uno de los \texttt{case
        \dots\_PAGE\_FAULT} (\texttt{trap.c:154-162}):
\begin{lstlisting}[style=cstyle]
case LOAD_PAGE_FAULT:
case LOAD_ACCESS_FAULT:
case INSTRUCTION_PAGE_FAULT:
case STORE_ACCESS_FAULT:
case STORE_PAGE_FAULT:
    printf("Task %s in CPU %d page fault. Killing task...\n", ...);
    task->exit_code = 1;
    task->killed = true;
    break;
\end{lstlisting}
    \item Luego del \texttt{switch}, si \texttt{task->killed},
        llama a \texttt{terminate()} (\texttt{task.c:237-259}) que
        pone \texttt{state=ZOMBIE}, hace \texttt{wakeup(t)} y entra
        a \texttt{sched}.
    \item El scheduler, en su próxima iteración, ve el ZOMBIE y
        llama \texttt{free\_resources(t)} (\texttt{task.c:82-87})
        que libera el kstack y pone \texttt{state=UNUSED}.
\end{enumerate}

\textbf{Nota:} EDOS \emph{no} implementa demand paging, COW ni
expansión del stack. Cualquier page fault de usuario es fatal. Es
una diferencia importante con xv6 (que tampoco hace demand paging
para texto, pero sí mata el proceso de forma similar; eso sí, en
xv6 más recientes el stack crece con guard).
\end{respuesta}

\subsection*{P9. Caminá un \texttt{sleep(4)}.}

\begin{respuesta}
\begin{enumerate}[leftmargin=*,nosep]
    \item User-space: \texttt{sleep(4)} (\texttt{usys.s:42-45})
        $\to$ \texttt{li a7, 5; ecall}.
    \item Trap a \texttt{u\_trap} $\to$ \texttt{user\_trap} $\to$
        \texttt{syscall} $\to$ \texttt{sys\_sleep}
        (\texttt{syscall.c:48-54}): \texttt{sleep(4)} en kernel.
    \item \texttt{sleep(4)} (\texttt{task.c:203-212}):
        \texttt{acquire(\&ticks\_lock)}, \texttt{task->sleep\_ticks
        = 4}, \texttt{suspend(\&ticks, \&ticks\_lock)}.
    \item \texttt{suspend} libera \texttt{ticks\_lock}, marca
        \texttt{state=WAITING, wait\_condition=\&ticks}, hace
        \texttt{sched(task)} — la tarea cede la CPU.
    \item Cada timer en el hart 0 llama a \texttt{inc\_ticks()}
        (\texttt{task.c:272-288}), que recorre las tareas
        \texttt{WAITING} con \texttt{wait\_condition == \&ticks} y
        decrementa su \texttt{sleep\_ticks}. Cuando llega a 0, pone
        \texttt{state=RUNNABLE}.
    \item En la siguiente iteración del scheduler, la tarea es
        elegible; \texttt{context\_switch} vuelve a \texttt{sched}
        en \texttt{suspend}, que reasume \texttt{ticks\_lock},
        retorna, y \texttt{sys\_sleep} devuelve 0.
    \item \texttt{user\_trap} retorna; \texttt{u\_trap\_ret} salta
        de vuelta a \texttt{ret} en \texttt{usys.s}.
\end{enumerate}
\end{respuesta}

\subsection*{P10. Caminá el camino del timer en kernel mode.}

\begin{respuesta}
\begin{enumerate}[leftmargin=*,nosep]
    \item Estamos ejecutando código del kernel (por ejemplo, dentro
        del scheduler o de una syscall).
    \item \texttt{mtime} alcanza \texttt{mtimecmp}; el hardware
        levanta IRQ M-mode.
    \item \texttt{mtvec = m\_trap}; saltamos a \texttt{m\_trap}
        (\texttt{trap.s:198-250}): salva los \texttt{a*}, llama a
        \texttt{next\_timer\_interrupt(tp)} para reprogramar
        \texttt{mtimecmp}, escribe \texttt{sip = 2} (\texttt{SSIP}),
        \texttt{mret}.
    \item De vuelta en S-mode, vemos pendiente la software int.
        \texttt{scause = 0x80000001} (TIMER\_INTERRUPT).
    \item Como \texttt{stvec = s\_trap}, saltamos a
        \texttt{s\_trap} (\texttt{trap.s:10}). Guarda registros en
        \emph{el mismo stack} (estamos en kernel), llama
        \texttt{kernel\_trap}.
    \item \texttt{kernel\_trap} (\texttt{trap.c:29-67}): en
        \texttt{case TIMER\_INTERRUPT}, incrementa ticks (sólo
        hart 0), \texttt{ack\_timer\_interrupt()} (limpia
        \texttt{sip.SSIP}); si hay tarea actual y se le acabó el
        quantum, \texttt{yield()}.
    \item Después del switch, \texttt{kernel\_trap} restaura
        \texttt{sepc} y \texttt{sstatus} (puede haber cambiado por
        un \texttt{context\_switch}) y retorna.
    \item \texttt{s\_trap} restaura registros y hace \texttt{sret}
        $\to$ vuelve al PC que tenía antes.
\end{enumerate}
\end{respuesta}

\section{¿Qué pasa si…?}

\subsection*{P11. ¿Qué pasa si saco el primer \texttt{sfence.vma} de \texttt{set\_page\_table}?}

\begin{respuesta}
La TLB de la CPU sigue conteniendo entradas del \texttt{satp}
\emph{anterior}. Cuando hagamos accesos a memoria justo después de
\texttt{w\_satp(\dots)}, la MMU resolverá esas direcciones usando
las traducciones viejas, no las nuevas. En la práctica:

\begin{itemize}[nosep,leftmargin=*]
    \item Si pasamos del \texttt{pgtbl} de un proceso A al de B,
        accesos a las primeras instrucciones del kernel podrían
        traducirse contra la tabla de A (que está pendiente de
        flush). El resultado típico es leer instrucciones
        equivocadas o una excepción de violación de permisos.
    \item En SMP, peor: cada CPU mantiene su propia TLB; si A se
        movió de CPU, la TLB de la nueva CPU puede no estar
        sincronizada.
\end{itemize}

Es justamente lo que \texttt{sfence.vma} (sin argumentos) hace:
invalida \emph{todas} las traducciones cacheadas en la TLB local.
\end{respuesta}

\subsection*{P12. ¿Y si saco el segundo \texttt{sfence.vma}?}

\begin{respuesta}
Generalmente menos grave que sacar el primero, pero igualmente
incorrecto. El segundo asegura que, después de que escribimos en
una tabla de páginas (creando una nueva entrada), la próxima
referencia a esa VA vea la nueva entrada en vez de un viejo cache
miss/parcial.

Sin el segundo, podríamos hacer \texttt{map\_page} y luego acceder
inmediatamente a la VA mapeada y obtener un page fault espurio,
porque la TLB aún no llenó la entrada.
\end{respuesta}

\subsection*{P13. ¿Qué pasa si no guardo \texttt{sepc} antes de llamar \texttt{user\_trap}?}

\begin{respuesta}
\texttt{sepc} es un CSR \emph{compartido} para el modo S. Si
durante \texttt{user\_trap} ocurre otro trap (por ejemplo, una
syscall del propio kernel — no, eso no pasa, pero un page fault
del kernel sí), \texttt{sepc} se sobrescribe con el nuevo PC. Al
volver, \texttt{sret} saltaría al lugar equivocado.

El kernel necesita ``apropiarse'' del \texttt{sepc} actual y
guardarlo. En \texttt{u\_trap} (\texttt{trap.s:136-137}) eso se
hace antes de la llamada:

\begin{lstlisting}[style=asmstyle]
csrr    a0, sepc
sw      a0, 30 * 4 (sp)     # tf->pc
\end{lstlisting}

Y al retornar, \texttt{u\_trap\_ret} (\texttt{trap.s:156-157})
restaura:

\begin{lstlisting}[style=asmstyle]
lw      a0, 30 * 4 (sp)
csrw    sepc, a0
\end{lstlisting}

Si lo olvidamos, después de un timer en U-mode el \texttt{sret}
final saltaría a la dirección del último trap que haya ocurrido.
\end{respuesta}

\subsection*{P14. ¿Y si el spinlock no deshabilita interrupciones?}

\begin{respuesta}
Deadlock clásico de tipo \emph{lock-to-self}:

\begin{enumerate}[leftmargin=*,nosep]
    \item Llamamos \texttt{acquire(\&L)}. Lo obtenemos.
    \item Llega un timer. \texttt{s\_trap $\to$ kernel\_trap $\to$
        yield}. \texttt{yield} a su vez quiere usar
        \texttt{tasks[i].lock} para cambiar de estado.
    \item Si el handler del trap necesita el lock \texttt{L} (o
        cualquier lock que esté actualmente tomado), giraría
        esperándolo \emph{en la misma CPU que ya lo tiene}.
        \textbf{Deadlock irrecuperable.}
\end{enumerate}

Aún cuando los locks no se solapen, hay un escenario más sutil:
si un \texttt{wakeup} se dispara desde un handler de trap mientras
nosotros tenemos el lock de \texttt{tasks[i]}, y el \texttt{wakeup}
intenta tomar ese mismo lock, esperamos eternamente.

Por eso \texttt{acquire} llama a \texttt{push\_irq\_off} antes del
test-and-set (\texttt{spinlock.c:29-30}).
\end{respuesta}

\subsection*{P15. ¿Qué pasa si \texttt{exec} no copia el kernel pgtbl?}

\begin{respuesta}
El proceso quedaría con un \texttt{pgtbl} que mapea su código y
stack, pero \emph{nada} del kernel. Apenas saliera por
\texttt{sret} a \texttt{u\_trap\_ret} (que está en \texttt{.text}
del kernel), generaría un \texttt{INSTRUCTION\_PAGE\_FAULT} porque
el PC apunta a una VA no mapeada en su tabla.

Lo más temprano que se vería es al ocurrir el primer trap: la CPU
salta a \texttt{stvec=u\_trap}. Esa dirección sería inválida en el
\texttt{pgtbl} actual del proceso. Result\-ado: una \emph{trap
loop} infinita o una excepción doble que termina en hardware reset.

La copia \texttt{memcpy(new\_pgtbl, kernel\_pgtbl, PAGE\_SIZE)} en
\texttt{task.c:386} es lo que evita esto.
\end{respuesta}

\subsection*{P16. ¿Qué pasaría si \texttt{NCPU} no coincide con \texttt{-smp N} de QEMU?}

\begin{respuesta}
\texttt{NCPU} dimensiona \texttt{cpus\_state[NCPU]}
(\texttt{task.c:14}). Si QEMU lanza más harts que \texttt{NCPU},
\texttt{cpuid()} devolvería un índice fuera de rango y todas las
operaciones sobre \texttt{cpus\_state[cpu]} corromperían memoria.

Si QEMU lanza menos, no pasa nada: los slots extra de
\texttt{cpus\_state[]} quedan sin usar.

En el \texttt{Makefile} actual se lanza \texttt{-smp 1}, así que
sólo el hart 0 corre y el resto del array está ocioso.
\end{respuesta}

\section{Cálculos sobre paginación Sv32}

\subsection*{P17. ¿Cuántas entradas hay por tabla en Sv32? ¿Cuánto ocupa cada tabla?}

\begin{respuesta}
Cada tabla (directorio raíz o tabla hoja) tiene \textbf{$2^{10}=1024$
entradas}. Cada entrada (PTE) es de \textbf{4 bytes}. Por lo tanto
una tabla ocupa exactamente \textbf{4096 B = 1 página}. Es por eso
que \texttt{alloc\_page} alcanza para crear un directorio o una
tabla hoja: una llamada $=$ una tabla.
\end{respuesta}

\subsection*{P18. Traducir manualmente la VA \texttt{0x80004ABC} con \texttt{kernel\_pgtbl}.}

\begin{respuesta}
\begin{itemize}[leftmargin=*,nosep]
    \item \texttt{VPN[1]} = bits 31..22 de \texttt{0x80004ABC} =
        \texttt{0b10\_0000\_0000} = \texttt{0x200} = 512.
    \item \texttt{VPN[0]} = bits 21..12 = \texttt{0b00\_0000\_0100}
        = \texttt{0x4} = 4.
    \item Offset = bits 11..0 = \texttt{0xABC} = 2748.
\end{itemize}

Buscamos \texttt{kernel\_pgtbl[512]}. \texttt{map\_kernel\_memory}
mapeó \texttt{ktext\_start = 0x80000000} con tamaño =
\texttt{\_\_text\_end - \_\_kernel\_start}, R+X, 1:1. Como
$0x80004ABC \in [0x80000000, \texttt{\_\_text\_end})$, \emph{si
\texttt{\_\_text\_end} > 0x80005000}, la VA está mapeada con
PTE para frame físico \texttt{0x80004} (es decir
\texttt{PPN=0x80004}, escrito \texttt{0x80004 << 10 | (R|X|V) =
0x20001000 | 0xb} \dots).

Bajando dos niveles:

\begin{itemize}[nosep,leftmargin=*]
    \item \texttt{root\_table[512]} apunta a la tabla hoja con su
        \texttt{PPN} y \texttt{V=1}.
    \item En esa tabla hoja, la entrada $[4]$ tiene \texttt{PPN =
        0x80004}, flags R+X+V.
    \item PA $= (\texttt{PPN} \ll 12) + \texttt{offset} = 0x80004000
        + 0xABC = $ \textbf{\texttt{0x80004ABC}}.
\end{itemize}

Como esperábamos para una región 1:1.
\end{respuesta}

\subsection*{P19. ¿Cuánta memoria consume el \emph{conjunto} de tablas de páginas de \texttt{init} en EDOS?}

\begin{respuesta}
Estimemos suponiendo que el \texttt{init.bin} pesa $\sim 1\,300$ B
(redondeamos a 1 página de código/datos) y la stack ocupa 1 página
más en \texttt{PROC\_MAX\_VA-4\,KiB}.

\begin{itemize}[leftmargin=*,nosep]
    \item Directorio raíz: 1 página (\texttt{exec} la asigna en
        \texttt{task.c:381}).
    \item Tabla hoja para \texttt{VPN[1]=0} (cubre la VA 0): 1
        página (la crea \texttt{get\_pte} cuando se mapea VA=0).
    \item Tabla hoja para \texttt{VPN[1]=511} (cubre la VA
        \texttt{0x7FC00000..0x80000000}, donde está el stack en
        \texttt{0x7FFFF000}): 1 página.
\end{itemize}

\textbf{Total para el proceso: 3 páginas $= $ 12\,KiB} (más las
PTs heredadas del kernel — \emph{compartidas} con todos los
procesos, no se cuentan).

Si el init fuera más grande y cruzara una frontera de 4\,MiB
(\texttt{VPN[1]} distinto), necesitaríamos más tablas hoja.
\end{respuesta}

\subsection*{P20. ¿Cuál es el máximo espacio de direcciones de un proceso EDOS?}

\begin{respuesta}
\texttt{PROC\_MAX\_VA = 0x80000000} = 2\,GiB (\texttt{task.h:20}).
Eso es de 0 a \texttt{0x7FFFFFFF} (excluido \texttt{0x80000000}).

¿Por qué no los 4\,GiB completos del Sv32? Porque desde
\texttt{0x80000000} hacia arriba está mapeada la RAM del kernel
1:1 en \texttt{kernel\_pgtbl}, y \texttt{exec} copia esa porción
\emph{en cada proceso}. Si dejáramos al usuario usar esas VAs,
chocaría con código del kernel y la separación user/kernel
perdería sentido.

Es exactamente la decisión que xv6 toma con \texttt{MAXVA} y
\texttt{KERNBASE}, aunque allí los números son distintos (Sv39).
\end{respuesta}

\section{Trade-offs}

\subsection*{P21. ¿Por qué Sv32 y no Sv39?}

\begin{respuesta}
La cátedra ya tiene que enseñar paginación, interrupciones, traps,
syscalls, planificación y filesystem en un cuatrimestre. Sv39
agrega:

\begin{itemize}[nosep,leftmargin=*]
    \item Tercer nivel (PD $\to$ PT inter. $\to$ PT hoja).
    \item Páginas de \textbf{4\,KiB, 2\,MiB y 1\,GiB} (megapages
        y gigapages) que complican el formato de PTE.
    \item Direcciones virtuales con \emph{sign extension} a 64 b
        (la regla de bits 63..39 = bit 38).
\end{itemize}

Ninguno de esos detalles es central para entender ``cómo funciona
la traducción de VA''. Sv32 da exactamente lo mismo
conceptualmente con la mitad del esfuerzo cognitivo.
\end{respuesta}

\subsection*{P22. ¿Por qué round-robin y no MLFQ o prioridades?}

\begin{respuesta}
Por la misma razón que Sv32: simplicidad. Un MLFQ requiere:

\begin{itemize}[nosep,leftmargin=*]
    \item Estructuras de colas múltiples.
    \item Reglas de subida/bajada de prioridad.
    \item Tracking del tiempo de CPU por tarea (más allá del cuanto).
\end{itemize}

Para que el alumno entienda ``cambio de contexto'',
``preemption'', ``estado del PCB'' y ``yield'', round-robin
alcanza. Las complicaciones de planificación se ven en teoría y
en los prácticos de scheduling (Práctico 4), no en EDOS.
\end{respuesta}

\subsection*{P23. ¿Por qué EFS y no xv6fs?}

\begin{respuesta}
Tres razones:

\begin{enumerate}[nosep,leftmargin=*]
    \item \textbf{Boot rápido}: no hay que montar un disco virtio,
        leer superblock, parsear inodes\dots\ basta \texttt{xxd
        -i} y un \texttt{efs\_files\_table[]}.
    \item \textbf{Foco}: el contenido conceptual de ``filesystem''
        en EDOS es ``mapeo nombre $\to$ bytes''. EFS lo cubre. Los
        detalles de journaling, FAT, inodes, etc., se cubren en
        prácticos teóricos.
    \item \textbf{Reproducibilidad}: como el FS vive en el binario,
        no hay que distribuir una imagen de disco aparte.
\end{enumerate}
\end{respuesta}

\subsection*{P24. ¿Por qué no hay \texttt{fork} en EDOS?}

\begin{respuesta}
Porque \texttt{fork} requiere \textbf{copiar el espacio de
direcciones del padre} (o implementar COW). Eso significa caminar
la page table del padre y, para cada PTE válida con \texttt{U=1},
o bien:

\begin{itemize}[nosep,leftmargin=*]
    \item Asignar una página nueva, hacer \texttt{memcpy}, mapearla
        en el pgtbl del hijo, \textbf{o}
    \item Marcar la PTE como solo-lectura en ambos pgtbls y manejar
        el page fault de COW.
\end{itemize}

Ambas opciones son ``ejercicios para el lector'' que pueden
agregarse en talleres posteriores. La versión 08-process se
contenta con \texttt{create\_process(path)}, que es esencialmente
\texttt{fork + exec} fusionados.
\end{respuesta}

\section{Comparativas con xv6}

\subsection*{P25. ¿Cómo se diferencia el manejo de la TLB?}

\begin{respuesta}
Ambos hacen \texttt{sfence.vma} antes y después de cambiar
\texttt{satp}. La diferencia es que xv6 tiene una rutina explícita
\texttt{kvminithart} llamada \emph{en cada hart después de boot}
para cargar \texttt{satp} apuntando al \texttt{kernel\_pgtbl}; en
EDOS el equivalente es \texttt{enable\_paging()} en
\texttt{kmain.c:28}, también llamado por cada hart antes de entrar
al \texttt{scheduler}.

xv6 mantiene además la VA \texttt{TRAMPOLINE} mapeada en todos los
pgtbls, lo que le permite hacer \texttt{satp} switches \emph{en
ensamblador} sin perder el código que ejecuta. EDOS evita esto
copiando el directorio del kernel en cada pgtbl de proceso.
\end{respuesta}

\subsection*{P26. ¿Qué hace xv6 con el trapframe que EDOS hace distinto?}

\begin{respuesta}
xv6:

\begin{itemize}[nosep,leftmargin=*]
    \item Reserva una página \emph{específica} para el trapframe de
        cada proceso, mapeada en una VA fija (\texttt{TRAPFRAME =
        TRAMPOLINE - PGSIZE}).
    \item \texttt{trampoline.S::uservec} accede a esa VA para
        guardar registros, sin tocar el stack del proceso.
    \item Después de eso, cambia \texttt{satp} al kernel pgtbl y
        salta al handler en C.
\end{itemize}

EDOS:

\begin{itemize}[nosep,leftmargin=*]
    \item Construye el trapframe \emph{en la cima del kernel-stack}
        del task. Su dirección es \texttt{kstack + 4K -
        sizeof(tf)}.
    \item No cambia \texttt{satp} en \texttt{u\_trap}; sigue con el
        pgtbl del proceso (que tiene los mapeos del kernel
        copiados).
    \item Cambia \texttt{satp} dentro de \texttt{user\_trap}
        (\texttt{trap.c:126}) sólo cuando ya está manejando la
        causa.
\end{itemize}

La razón didáctica: EDOS prefiere ``todo está en la pila'' a
``buscá la página mágica''.
\end{respuesta}

\subsection*{P27. ¿Cómo se diferencia el sleep-lock de xv6 con \texttt{suspend} en EDOS?}

\begin{respuesta}
xv6 tiene \textbf{dos} tipos de locks:

\begin{itemize}[nosep,leftmargin=*]
    \item \texttt{struct spinlock}: el clásico, no duerme.
    \item \texttt{struct sleeplock}: encima de \texttt{spinlock},
        con \texttt{acquiresleep/releasesleep}. Si no se puede
        tomar, llama a \texttt{sleep(\&slk->locked, \&slk->lk)},
        que pone al proceso \texttt{SLEEPING} hasta que el dueño lo
        libere y llame \texttt{wakeup}.
\end{itemize}

EDOS tiene \emph{sólo} \texttt{spinlock} más
\texttt{suspend/wakeup} (\texttt{task.c:163,186}). Lo segundo es
construible \emph{encima de} lo primero — y de hecho la API
\texttt{suspend(condition, lk)} es funcionalmente equivalente a un
sleep-lock que duerme en \texttt{condition} y reasume el lock al
volver. EDOS no encapsula el patrón en un tipo separado.
\end{respuesta}

\subsection*{P28. ¿Qué syscalls tiene xv6 que no tiene EDOS?}

\begin{respuesta}
xv6 expone alrededor de 22 syscalls; EDOS, 6. Las que faltan en
EDOS:

\begin{itemize}[nosep,leftmargin=*]
    \item \texttt{fork}, \texttt{wait}: no hay relación padre-hijo
        real en EDOS.
    \item \texttt{exec}: existe en kernel (\texttt{task.c:376})
        pero no como syscall.
    \item \texttt{read}, \texttt{write}, \texttt{open},
        \texttt{close}, \texttt{dup}, \texttt{fstat}: no hay
        descriptores de archivo en EDOS.
    \item \texttt{pipe}: no hay pipes.
    \item \texttt{kill}: no hay señales.
    \item \texttt{mkdir}, \texttt{link}, \texttt{unlink},
        \texttt{chdir}: el FS de EDOS no soporta jerarquía ni
        escrituras.
    \item \texttt{sbrk}: no hay heap manage\-able en EDOS (el
        stack es de tamaño fijo, una página).
\end{itemize}

A cambio EDOS tiene unas pocas convenientes (\texttt{console\_puts}/
\texttt{console\_putc} explícitas), que xv6 implementa vía
\texttt{write(1, \dots)}.
\end{respuesta}

\section{Detalles finos del código}

\subsection*{P29. ¿Por qué \texttt{m\_trap} pone los \texttt{a*} en 0 después de salvarlos?}

\begin{respuesta}
Mirá \texttt{trap.s:224-239}:

\begin{lstlisting}[style=asmstyle]
lw      a0, 0 * 4 (sp)
sw      x0, 0 * 4 (sp)
\dots
\end{lstlisting}

Restaura y \emph{escribe cero} sobre el slot del stack que acaba
de leer. ¿Para qué? Para que, si después el código de un proceso
de usuario lee el stack del kernel (por algún error del SO o por
una vulnerabilidad de filtración), no encuentre valores
identificables. Es una mitigación defensiva preventiva, mínima
pero presente.
\end{respuesta}

\subsection*{P30. ¿Por qué \texttt{u\_trap} hace \texttt{csrrw sp, sscratch, sp}?}

\begin{respuesta}
\texttt{csrrw rd, csr, rs1} es ``intercambio atómico'': escribe
\texttt{rs1} en \texttt{csr} y deja en \texttt{rd} el valor previo.
La instrucción aquí cambia \texttt{sp} (que tenía el SP de usuario)
por \texttt{sscratch} (que tenía la base del kstack del proceso),
en \textbf{una sola instrucción}. Esto es importante porque, al
entrar al trap, todavía estamos con stack de usuario; cualquier
escritura previa al swap rompería el stack del proceso.

Después del swap, \texttt{sp} apunta al kstack y \texttt{sscratch}
al user-stack. La siguiente instrucción reserva espacio en kstack
para el trap frame, y unas líneas más abajo
(\texttt{trap.s:132-133}) \texttt{csrr a0, sscratch; sw a0, 29*4(sp)}
guarda el SP de usuario en el frame.
\end{respuesta}

\subsection*{P31. ¿Por qué \texttt{create\_task} hace \texttt{memset(\&ctx, 0, sizeof ctx)} antes de asignar \texttt{ra}?}

\begin{respuesta}
Para que los registros \texttt{s0--s11} y \texttt{sp} en el
contexto inicial estén en 0. Cuando \texttt{context\_switch} los
cargue al programar la primera ejecución, esos \texttt{s*}
quedarán en cero (que es lo que correspondería para una función
recién entrada). Sin el \texttt{memset}, leeríamos basura del slot
\texttt{tasks[i]} que tal vez quedó de una tarea previa
\texttt{UNUSED}.
\end{respuesta}

\subsection*{P32. ¿Por qué \texttt{exec} libera la página vieja \emph{después} de mapear la nueva?}

\begin{respuesta}
Por seguridad: si la asignación o el \texttt{load\_program} fallan
(\texttt{goto error}), todavía conservamos la tabla del proceso
original y podemos devolver \texttt{false} sin haberlo dejado en un
estado inválido.

La lógica está en \texttt{task.c:406-411}:

\begin{lstlisting}[style=cstyle]
if (is_process(task)) {
    unmap_region(task->pgtbl, 0, PROC_MAX_VA, true);
    free_page(task->pgtbl);
}
task->pgtbl = new_pgtbl;
\end{lstlisting}

Sólo después de que \texttt{new\_pgtbl} está completo y validado
liberamos el viejo. Es el patrón \emph{prepare-then-swap} típico.
\end{respuesta}

\subsection*{P33. ¿Cuándo se incrementa \texttt{ticks} y por qué sólo el hart 0?}

\begin{respuesta}
En \texttt{trap.c:46-48} (en \texttt{kernel\_trap}) y
\texttt{trap.c:133-135} (en \texttt{user\_trap}), dentro del caso
\texttt{TIMER\_INTERRUPT}, se hace:

\begin{lstlisting}[style=cstyle]
if (cpu_id == 0)
    inc_ticks();
\end{lstlisting}

Cada hart tiene su propio \texttt{mtimecmp}, así que recibirían su
propio timer y cada uno incrementaría ticks, llevando la cuenta al
\emph{múltiplo del número de harts}. Para mantener ``ticks $=$
tiempo real'', sólo el hart 0 incrementa.

Es una decisión simple: en xv6 hay una similar (sólo el hart 0
mantiene la noción de tiempo).
\end{respuesta}

\subsection*{P34. ¿Qué hace \texttt{set\_u\_previous\_mode} y por qué hace falta?}

\begin{respuesta}
\texttt{arch.h:411-417}:

\begin{lstlisting}[style=cstyle]
static inline void set_u_previous_mode(void) {
    size_t status = cpu_status();
    status &= ~(1 << 8);    // clear sstatus.SPP
    status |= (1 << 5);     // sstatus.SPIE = 1
    set_cpu_status(status);
}
\end{lstlisting}

\texttt{sstatus.SPP} (Supervisor Previous Privilege) en 0 indica
``el modo previo era U'', así que el próximo \texttt{sret} bajará a
U-mode. \texttt{sstatus.SPIE} (Previous Interrupt Enable) en 1
indica ``cuando se ejecute \texttt{sret}, hay que dejar
\texttt{SIE=1}'' (es decir, interrupciones habilitadas en U-mode).

Sin esta función, después del primer \texttt{return\_to\_user\_mode}
\texttt{sret} no haría el cambio de modo correcto.
\end{respuesta}

\subsection*{P35. ¿Por qué \texttt{sleep} suspende sobre \texttt{\&ticks} y no sobre el propio task?}

\begin{respuesta}
Porque la condición no es ``este task termine'', sino ``\emph{los
ticks alcancen un valor}''. Toda tarea durmiendo por tiempo
comparte la misma fuente de wakeups: el timer en \texttt{inc\_ticks}.
\texttt{inc\_ticks} (\texttt{task.c:272-288}) hace un solo recorrido
de la tabla, y para todas las tareas con
\texttt{wait\_condition == \&ticks} decrementa su
\texttt{sleep\_ticks}.

Si se durmiera sobre un canal por tarea, \texttt{inc\_ticks}
tendría que enviar wakeups individuales, lo cual es más caro y no
agrega claridad.
\end{respuesta}

% =====================================================================
\part{Cheat sheet final}
% =====================================================================

\section{Registros RISC-V relevantes (1 cara)}

\begin{center}
\footnotesize
\begin{tabular}{@{}llp{9.5cm}@{}}
\toprule
Reg./CSR & Modo & Rol \\
\midrule
\texttt{x0}             & --- & Siempre cero (no se puede escribir). \\
\texttt{ra}             & --- & Return address. Es el primer registro guardado por \texttt{context\_switch}. \\
\texttt{sp}             & --- & Stack pointer. Cada hart arranca con un stack propio en \texttt{\_\_stack0}. \\
\texttt{gp}             & --- & Global pointer (set por el linker). \\
\texttt{tp}             & --- & \emph{En EDOS, el hartid}. Se setea en \texttt{arch.s:14}. \\
\texttt{a0--a7}         & --- & Argumentos y resultado. \texttt{a0} también porta el retorno de syscalls. \\
\texttt{s0--s11}        & --- & Callee-saved. Es lo que \texttt{context\_switch} salva además de \texttt{ra} y \texttt{sp}. \\
\texttt{t0--t6}         & --- & Caller-saved. Sólo se salvan en el trap frame. \\
\midrule
\texttt{mhartid}        & M   & Identifica el hart. Leído en boot. \\
\texttt{mstatus}        & M   & \texttt{MPP} (modo previo), \texttt{MIE} (habilita IRQs M). \\
\texttt{mepc}           & M   & PC para \texttt{mret}. \\
\texttt{mtvec}          & M   & Handler M-mode (en EDOS, \texttt{m\_trap}). \\
\texttt{mie}            & M   & Bit \texttt{MTIE} habilita IRQ del timer. \\
\texttt{medeleg / mideleg} & M & Delega exc./int.\ a S. \\
\midrule
\texttt{sstatus}        & S   & \texttt{SIE} (habilita IRQs), \texttt{SPP} (modo previo: 0=U, 1=S), \texttt{SPIE} (estado previo de SIE). \\
\texttt{sepc}           & S   & PC en el trap, restaurado por \texttt{sret}. \\
\texttt{scause}         & S   & Causa del trap (bit 31 = interrupción). \\
\texttt{stval}          & S   & Dirección que causó el fault (page faults). \\
\texttt{stvec}          & S   & Handler S-mode (en EDOS, \texttt{s\_trap} o \texttt{u\_trap}). \\
\texttt{sscratch}       & S   & ``Bolsillo'' del kernel. En EDOS: base + 4\,KiB del kstack del proceso. \\
\texttt{satp}           & S   & Page table register. \texttt{0x80000000 | PPN(root)} en Sv32. \\
\texttt{sie / sip}      & S   & Habilitación / pendiente de IRQs S-mode. \\
\midrule
\texttt{pmpcfg0, pmpaddr0} & M & Configuración del Physical Memory Protection. \\
\texttt{mtime / mtimecmp[h]} & M & Timer (MMIO en CLINT). \\
\bottomrule
\end{tabular}
\end{center}

\section{Formato de PTE Sv32}

\begin{center}
\begin{tikzpicture}[
    every node/.style={font=\ttfamily\small},
    cell/.style={draw,minimum height=0.7cm,align=center}
]
\node[cell,minimum width=5cm,fill=c1!20] (ppn) at (0,0) {PPN (22 bits)};
\node[cell,minimum width=0.6cm,fill=cfree!20,right=0pt of ppn] (r0) {res};
\node[cell,minimum width=0.5cm,fill=c3!25,right=0pt of r0]   (d)  {D};
\node[cell,minimum width=0.5cm,fill=c3!25,right=0pt of d]    (a)  {A};
\node[cell,minimum width=0.5cm,fill=c3!25,right=0pt of a]    (g)  {G};
\node[cell,minimum width=0.5cm,fill=c3!25,right=0pt of g]    (u)  {U};
\node[cell,minimum width=0.5cm,fill=c3!25,right=0pt of u]    (x)  {X};
\node[cell,minimum width=0.5cm,fill=c3!25,right=0pt of x]    (w)  {W};
\node[cell,minimum width=0.5cm,fill=c3!25,right=0pt of w]    (r)  {R};
\node[cell,minimum width=0.5cm,fill=c3!25,right=0pt of r]    (v)  {V};
\node[font=\scriptsize] at ($(ppn.north)+(0,0.25)$) {bits 31..10};
\node[font=\scriptsize] at ($(v.north)+(0,0.25)$) {0};
\end{tikzpicture}
\end{center}

\textbf{EDOS sólo usa V, R, W, X, U} (\texttt{arch.h:175-179}). En
Sv32 ``oficial'' hay además A (accessed), D (dirty) y G (global)
que el hardware actualiza, pero QEMU no los gestiona y EDOS no los
consulta.

\textbf{Reglas:} \texttt{V} debe estar siempre en 1 para una entrada
viva. Si \texttt{R|W|X = 0}, la entrada es un puntero a la próxima
tabla (no es hoja). Si alguno de R/W/X está prendido, es hoja.

\section{Secuencia mínima de boot}

\begin{enumerate}[leftmargin=*,nosep]
    \item \texttt{boot} en \texttt{arch.s}: PMP, MPP=S, deleg,
        \texttt{sie}, \texttt{mret} a \texttt{supervisor}.
    \item \texttt{supervisor:} en \texttt{arch.s}: \texttt{stvec=s\_trap},
        \texttt{kernel\_main}.
    \item \texttt{kernel\_main} en \texttt{kmain.c}: hart 0 hace
        \texttt{init\_kalloc, init\_vm, create\_process("init")},
        publica \texttt{ready=1}.
    \item Todos los harts: \texttt{enable\_paging, scheduler}.
\end{enumerate}

\section{Secuencia del primer scheduling de init}

\begin{enumerate}[leftmargin=*,nosep]
    \item Scheduler ve \texttt{tasks[k].state == RUNNABLE} con
        \texttt{ctx.ra=start\_process}.
    \item \texttt{context\_switch} $\to$ \texttt{ret} a
        \texttt{start\_process}.
    \item \texttt{start\_process}: \texttt{release(\&t->lock); return\_to\_user\_mode();}
    \item \texttt{return\_to\_user\_mode}: \texttt{stvec=u\_trap,
        sscratch=kstack+4K, SPP=0, SPIE=1, satp=task->pgtbl,
        u\_trap\_ret(tf);}
    \item \texttt{u\_trap\_ret}: \texttt{sp=tf, sepc=tf->pc=0,
        restaurar regs, sret}.
    \item \texttt{sret} $\to$ U-mode @ \texttt{pc=0} (entry
        \texttt{start} de \texttt{init}).
\end{enumerate}

\section{Secuencia de una syscall}

\begin{enumerate}[leftmargin=*,nosep]
    \item User: \texttt{li a7, N; ecall}.
    \item HW $\to$ \texttt{u\_trap}: swap \texttt{sp/sscratch},
        salvar regs, llamar \texttt{user\_trap}.
    \item \texttt{user\_trap}: \texttt{satp=kernel\_pgtbl,
        stvec=s\_trap}, dispatch por \texttt{scause=8},
        \texttt{syscall(task)}, \texttt{skip\_trap\_instruction}.
    \item \texttt{syscall}: lee \texttt{a7}, llama
        \texttt{syscalls\_table[N](task)}, escribe \texttt{tf->a0}.
    \item \texttt{return\_to\_user\_mode} $\to$ \texttt{u\_trap\_ret}
        $\to$ \texttt{sret}.
\end{enumerate}

\section{Comandos útiles para QEMU + GDB}

\begin{lstlisting}[style=shellstyle]
# Compilar y correr
make
make qemu          # arranca QEMU con -smp 1

# Salir de QEMU
Ctrl-A, x

# Ver el ensamblado del kernel
less kernel.asm

# Ver el ensamblado de init
less user/init.asm

# Debug con GDB:
# Terminal 1:
qemu-system-riscv32 -machine virt -bios none -nographic -smp 1 \
                    -kernel kernel -s -S
# (-s = gdbserver en :1234, -S = freezar al inicio)

# Terminal 2:
riscv64-unknown-elf-gdb kernel
(gdb) target remote :1234
(gdb) set arch riscv:rv32
(gdb) b kernel_main
(gdb) c

# Para ver registros CSR:
(gdb) info registers
(gdb) p/x $satp
(gdb) p/x $sepc

# Para imprimir memoria como un trap_frame:
(gdb) p *(struct trap_frame *)(task->kstack + 4096 - sizeof(struct trap_frame))

# Para ver dónde está mapeada una VA en el pgtbl del proceso:
(gdb) p/x ((unsigned int*)task->pgtbl)[0]            # PDE[0]
(gdb) p/x ((unsigned int*)pte_pa(PDE[0]))[0]         # PTE[0] de la hoja
\end{lstlisting}

\section{Tabla mental rápida}

\begin{center}
\small
\begin{tabular}{@{}lll@{}}
\toprule
Pregunta urgente & Mirar en\dots & Línea \\
\midrule
``Cómo se monta un PTE'' & \texttt{arch.c} & 69-73 \\
``Cómo entra al kernel'' & \texttt{arch.s} & 11-83 \\
``Qué hace el timer'' & \texttt{trap.s} + \texttt{trap.c} & 198 + 45 \\
``Cómo cede la CPU una tarea'' & \texttt{task.c} & 140-158 \\
``Cómo cambia entre tareas'' & \texttt{arch.s} & 90-122 \\
``Cómo se carga un programa'' & \texttt{task.c} & 376-421 \\
``Por qué hay un \texttt{sscratch} & \texttt{trap.s} + \texttt{trap.c} & 93 + 84 \\
``Cómo despierta sleep'' & \texttt{task.c} & 272-288 \\
``Dónde está la tabla de syscalls'' & \texttt{syscall.c} & 63-70 \\
``Cómo se inicializa la VM'' & \texttt{vm.c} + \texttt{arch.c} & 6 + 97 \\
\bottomrule
\end{tabular}
\end{center}

\end{document}
